\providecommand{\recsinc}{\operatorname{sinc}}
\providecommand{\recgE}{\gamma_{\mathrm{E}}}
\providecommand{\recGen}{\operatorname{Gen}}
\providecommand{\recav}{\operatorname{av}}
\providecommand{\recA}{\mathcal{A}}
\providecommand{\recartanh}{\operatorname{artanh}}
\section{Matching Euler factors of products of Dirichlet \texorpdfstring{$L$}{L}-functions}\label{sec:eul}

How often can two products of Dirichlet $L$-functions have the same Euler factor at a prime? For
characters of exponent two, of cyclic groups, of abelian $\ell$-groups and of odd prime order this
section determines, in the cases stated in each subsection, the least density of primes at which the
local factors differ, the extremal pairs, and the length of a prime window that exhibits a difference.
The only input from Dirichlet GRH is the prime-sum estimate \cref{rec:lem:gram} in windows (\stG); the rest is finite group theory, and no zero is located.
The local polynomials, the divisor criterion \cref{rec:prop:cdivisor} and the single prime-window
estimate \eqref{rec:eq:cgram} are stated once and used for all exponents. \Cref{rec:sec:cmatch} precedes
\cref{rec:sec:omatch}, whose proofs use \cref{rec:lem:pproj,rec:thm:pmass,rec:thm:psat}; at odd prime
order the saturated and prime-power minima and the coset equality case are quoted from there.

\subsection{Matching Euler factors of products of quadratic $L$-functions}\label{rec:sec:qmatch}
\sources{dirichlet_packet_reconstruction}

Let $X$ be a group of characters modulo $Q$ of order dividing two, $r=\dim_{\F_2}X\ge1$, and $V=\widehat X$,
written additively. A packet $A$ is a multiset of $k\ge1$ members of $X$, and at $p\nmid Q$ its local
polynomial $\prod_{\chi\in A}(1-\chi(p)T)$ is determined by the trace $\sum_{\chi\in A}\chi(p)$ (the case $E=2$ of
\cref{rec:prop:cdivisor}; compare \cref{rec:lem:otrace}). For distinct packets $A,C$ let
$\delta(A,C)$ be the density of primes with different local polynomials. By the prime number theorem in
progressions,
\[
 \delta(A,C)=2^{-r}\abs{\supp H},\qquad H(v)=\sum_{\chi\in X}c_\chi\chi(v),\quad c_\chi=\operatorname{mult}_A(\chi)-\operatorname{mult}_C(\chi),
\]
with $H\ne0$, $H(0)=0$, $B=\sum\abs{c_\chi}\le2k$, and after cancellation each packet has size $B/2$.

\begin{lemma}[Dyadic valuation; \stU]\label{rec:lem:walsh}
For a nonzero integral combination $H$ of Walsh characters of $\F_2^r$ and $t=\min_{H(x)\ne0}v_2(H(x))$, the map
$x\mapsto2^{-t}H(x)\bmod2$ is a nonzero Boolean polynomial of degree at most $t$. A nonzero Boolean
polynomial of degree at most $s\le r$ has relative weight at least $2^{-s}$, with equality only for the
indicator of an affine flat.
\end{lemma}

\begin{proof}
In Boolean coordinates each Walsh character is a product of factors $1-2x_i$, so a monomial of degree
$j$ has coefficient divisible by $2^j$, and by $2^t$ by M\"obius inversion; divide and reduce. For the
weight write $f_0+x_rf_1$ and induct. At equality, translates of the support coincide or are
disjoint, since a nonzero Boolean difference of degree at most $s-1$ has weight at least $2^{1-s}$.
\end{proof}

\begin{theorem}[Exact minimum and the high-agreement spectrum; \stU]\label{rec:thm:qmatch}
Put $D=\floor{\log_2(2k)}$ and $d=\min(r,D)$. The least positive disagreement density at rank $r$ and size
$k$ is $2^{-d}$. If $r>D$, equality holds exactly when $H=\pm2^d\eta\mathbf 1_{a+W}$ with $\eta\in X$ and an affine
subspace $a+W\not\ni0$ of codimension $d$; if $r\le D$, exactly when $H=m2^r\mathbf 1_{\set a}$, $a\ne0$, $\abs m2^r\le2k$.
If $r>D$, the densities strictly below $2^{1-d}$ are exactly
\[
 \delta_j=2^{1-d}(1-2^{-j}),\qquad1\le j\le\min(d,r-d),\quad k\ge2^d-2^{d-j},
\]
and every pair realizing $\delta_j$ reduces after cancellation to two disjoint sets of exactly $2^d-2^{d-j}$
distinct characters; if $r\le D$ only the minimum occurs below twice the minimum.
\end{theorem}

\begin{proof}
Since $\abs H\le2k$, $t\le D$, and \cref{rec:lem:walsh} gives $\delta\ge2^{-t}\ge2^{-d}$. At equality all nonzero values
are $\pm2^d$ on an affine flat $a+W$, $H(a+w)=2^ds(w)$, and Fourier inversion makes every normalized
coefficient of $s$ integral; by Parseval exactly one is $\pm1$, so $s$ is a signed character. The
converse uses the two halves of the Fourier support, padded by common entries. Below twice the
minimum $t=d$ and the Boolean polynomial is the disagreement indicator, whose weights are given by
the Kasami--Tokura classification of Reed--Muller weights below twice the minimum
\cite{KasamiTokura1970Weight}. Then $\abs{c_\chi}\le2^d\delta<2$, so $B=\sum c_\chi^2=4^d\delta$ and $B/2=2^d-2^{d-j}\le k$, and
$\abs{\supp H}$ is even, forcing $j\le r-d$. Two $d$-dimensional subspaces $U,U'\subset X$ meeting in dimension $d-j$
give $H=2^d(\mathbf 1_{U^\perp}-\mathbf 1_{(U')^\perp})$ with density $\delta_j$. For $r\le D$ densities are multiples of $2^{-r}$.
\end{proof}

For example, for $r\ge6$ two products of seven factors can agree at $25/32$ of the primes, but not at
any proportion strictly between $25/32$ and $13/16$. The classification is checked on $723{,}057$ pairs
and $95$ constructions through rank ten:
\begin{center}
\artifact{dirichlet_packet_reconstruction/prime-traces/verification.txt}
\end{center}

\begin{theorem}[Finite detection and density recovery; \stG]\label{rec:thm:qmatchfinite}
Let $M_x(f)=\sum_{p\nmid Q}(\log p)w(p/x)f([p])$ be the window sum of \eqref{rec:eq:cgram} for $f:V\to\C$, and $\norm f_A=\sum_\chi\abs{\widehat f(\chi)}$. If $r>D$ and
$0<\delta<2^{1-d}$, the disagreement set $S$ has $M_x(\mathbf 1_S)=x\delta+O_w(\sqrt x\log(2Q))$ uniformly in $k,r,Q$, and
rounding $4^dM_x(\mathbf 1_S)/(2x)$ recovers $B/2$, hence $\delta$, for $x\ge C_w'16^d\log^2(2Q)$.
\end{theorem}

\begin{proof}
In the high-agreement range $\mathbf 1_S=4^{-d}H^2$, so $\norm{\mathbf 1_S}_A<4$, and $4^d\delta/2=B/2$ by \cref{rec:thm:qmatch};
apply \eqref{rec:eq:cgram}.
\end{proof}

Every distinct pair has a distinguishing prime in $(x/2,x)$ once $x\ge C_wk^2\log^2(2Q)$, by \cref{rec:rem:cscale}
with $m=B/2\le k$.

\subsection{Matching Euler factors in cyclic groups and abelian $\ell$-groups}\label{rec:sec:cmatch}
\sources{euler_agreement_2026-11-27}

At exponent two (\cref{rec:sec:qmatch}) and at odd prime exponent (\cref{rec:sec:omatch}) a single trace decides
the local polynomial. Here $X$ is an arbitrary nontrivial finite abelian group of Dirichlet characters
modulo $Q$, written additively, $V=\Hom(X,\C^\times)$, and $v_p\in V$ is given by $v_p(x)=x(p)$ for $p\nmid Q$. A
packet of size $k$ is a function $A:X\to\Z_{\ge0}$ with $\sum_xA(x)=k$, with local polynomial
\[
 E_A(v,T)=\prod_{x\in X}(1-v(x)T)^{A(x)}.
\]
For distinct packets of equal size put $c=A-B$, $m=\sum_xc(x)^+=\sum_xc(x)^->0$ (the mass left after
removing common characters, so $m\le k$), $\mathcal D(c)=\set{v\in V:E_A(v,T)\ne E_B(v,T)}$,
\[
 \delta(c)=\frac{\abs{\mathcal D(c)}}{\abs X},\qquad d_X(k)=\min_{A\ne B,\ \abs A=\abs B=k}\delta(A-B),\qquad H_c(v)=\sum_xc(x)v(x).
\]

\begin{lemma}[Histograms, powers and density; \stU]\label{rec:lem:chist}
\begin{enumerate}
\item $v\notin\mathcal D(c)$ if and only if $\sum_{x:v(x)=z}c(x)=0$ for every root of unity $z$.
\item Agreement at $v$ implies agreement at $v^j$, $v^j(x)=v(x)^j$, for every integer $j$, and the two are
equivalent when $j$ is prime to the exponent of $X$.
\item The relative natural density of the primes $p\nmid Q$ at which the local polynomials differ exists and
equals $\delta(c)$.
\end{enumerate}
\end{lemma}

\begin{proof}
Distinct values $z$ give distinct linear factors $1-zT$, and unique factorization in $\C[T]$ gives (i).
Pushing equal value multisets through $z\mapsto z^j$ gives (ii), with the converse when this map permutes
the relevant roots of unity. For (iii), the evaluation map $(\Z/Q\Z)^\times\to V$ is onto, since a proper
image would have a nontrivial annihilator, a nontrivial member of $X$ equal to one at every unit; its
fibres have equal size, and the prime number theorem in progressions \cite{Davenport2000Multiplicative}
gives every reduced class the same density.
\end{proof}

\begin{proposition}[Divisor-indexed traces; \stU]\label{rec:prop:cdivisor}
If $E$ is the exponent of $X$, then $v\notin\mathcal D(c)$ if and only if $H_c(v^d)=0$ for every positive proper divisor
$d$ of $E$.
\end{proposition}

\begin{proof}
Fourier inversion of the signed histogram on the $E$th roots of unity says that it vanishes exactly when
all $H_c(v^j)$, $0\le j<E$, vanish. For $j=0$ the value is zero since the packets have equal size. For
$1\le j<E$ write $d=\gcd(j,E)$ and $j=du$; the automorphism of $\Q(\zeta_{E/d})$ taking $\zeta_{E/d}$ to $\zeta_{E/d}^u$ carries
$H_c(v^d)$ to $H_c(v^j)$, because $c$ is integral and $\gcd(u,E/d)=1$.
\end{proof}

One trace no longer suffices at order four: if $\chi(p)=i$, the packets $\set{1,\chi^2}$ and $\set{\chi,\chi^3}$ both have
trace zero at $p$, but their local polynomials are $1-T^2$ and $1+T^2$.

\paragraph{Cyclic groups.}
Let $X=\langle\chi\rangle$ be cyclic of order $E\ge2$, let $\ell_1<\dots<\ell_s$ be the distinct prime divisors of $E$, and put
\[
 t=t(E,k)=\min\set{s,\floor{\log_2(2k)}},\qquad D(E,k)=\prod_{i=1}^t\Bigl(1-\frac1{\ell_i}\Bigr).
\]

\begin{theorem}[Cyclic Euler matching; \stU]\label{rec:thm:cyclic}
For every cyclic character group $X$ of order $E\ge2$ and every $k\ge1$, $d_X(k)=D(E,k)$. The minimum depends
only on the distinct prime divisors of $E$ and on $k$, every dyadic plateau of the formula is attained
inside every such $X$, and the unrestricted minimum $\varphi(E)/E$ is first attained at $k=2^{s-1}$.
\end{theorem}

The finite input is a signed-array form of the classical minimum-volume principle for trades
\cite{GhorbaniKamaliKhosrovshahiKrotov2020Volumes}.

\begin{lemma}[A small nonzero marginal; \stU]\label{rec:lem:cmarginal}
Let $B_1,\dots,B_s$ be finite nonempty sets and $c:B_1\times\dots\times B_s\to\Z$ nonzero. For $I\subseteq\set{1,\dots,s}$ let
$c_I(x_I)=\sum_{x_{I^c}}c(x_I,x_{I^c})$ be its marginal ($c_\varnothing=\sum_xc(x)$). If $r=\min\set{\abs I:c_I\ne0}$, then $\norm c_1\ge2^r$.
Consequently, if $c_\varnothing=0$ and $\norm c_1=2m$, some nonempty $I$ with $\abs I\le\floor{\log_2(2m)}$ has $c_I\ne0$.
\end{lemma}

\begin{proof}
Induct on $s$. If $r=0$, integrality gives $\norm c_1\ge\abs{\sum c}\ge1$. Let $r\ge1$. As $c\ne0$ has total sum zero,
its support has two points; choose a coordinate $i$ in which two support points differ. There are at
least two nonzero slices $c^{(z)}$, obtained by fixing $x_i=z$. Every marginal of a slice on fewer than $r-1$ of
the remaining coordinates is an entry of a marginal $c_{I\cup\set i}$ on fewer than $r$ coordinates, hence zero.
By induction each nonzero slice has $\norm{c^{(z)}}_1\ge2^{r-1}$, and the slices are disjoint.
\end{proof}

\begin{proof}[Proof of \cref{rec:thm:cyclic}]
Identify $X$ and $V$ with $\Z/E\Z$, and write $c_x$ for the difference of multiplicities; after cancellation
$c\ne0$, $\sum_xc_x=0$ and $\norm c_1=2m\le2k$. By \cref{rec:lem:chist} the polynomials agree at $v$ exactly when the
signed histogram $(v_*c)(z)=\sum_{vx=z}c_x$ vanishes, and $\delta$ is the fraction of $v$ with $v_*c\ne0$.
Write $E=\prod_i\ell_i^{a_i}$ and use the coordinates $X\cong\prod_i\Z/\ell_i^{a_i}\Z$. By \cref{rec:lem:cmarginal} some nonzero
marginal $c_I$ uses $r=\abs I\le\min\set{s,\floor{\log_2(2m)}}\le t$ coordinates. If $\ell_i\nmid v_i$ for all $i\in I$, multiplication
by $v_i$ is bijective on these coordinates, so the $I$-marginal of $v_*c$ is an invertible relabelling of $c_I\ne0$.
Hence $\mathcal D(c)$ contains the cylinder $C_I=\set{v:\ell_i\nmid v_i\ (i\in I)}$, of density
\[
 d_I=\prod_{i\in I}(1-\ell_i^{-1})\ge\prod_{i=1}^r(1-\ell_i^{-1})\ge D(E,k).
\]
For equality let $h_i=E/\ell_i$ for $1\le i\le t$ and take the convolution $c=\mathop{*}_{i=1}^t(\delta_0-\delta_{h_i})$ in $\Z/E\Z$: $A$
consists of the sums $\sum_{i\in J}h_i$ over even $\abs J$, and $C$ of those over odd $\abs J$ (the \emph{parity box}). The
$2^t$ sums are distinct, since their projections to the primary components distinguish the choices
independently, and both packets have size $2^{t-1}\le k$. Under multiplication by $v$ a factor becomes zero
exactly when $\ell_i\mid v$; otherwise the same argument shows that the $2^t$ signed atoms remain distinct. So
$v_*c\ne0$ exactly on $C_{\set{1,\dots,t}}$, of density $D(E,k)$; pad both packets with the same $k-2^{t-1}$ characters.
The products $\prod_{i\le j}(1-\ell_i^{-1})$ decrease strictly in $j$, so $\varphi(E)/E$ is reached exactly when $t=s$, that is
$k\ge2^{s-1}$.
\end{proof}

Let $p_j$ be the $j$th prime ($p_1=2$) and $d(k)=\floor{\log_2(2k)}$.

\begin{corollary}[The global cyclic profile; \stU]\label{rec:cor:cglobal}
For every $k\ge1$, the infimum of $\delta$ over all cyclic groups of Dirichlet characters and all distinct
packets of size $k$ is $\prod_{j=1}^{d(k)}(1-p_j^{-1})$. It is a minimum, even when all factors are required to be
primitive of one prime conductor, and it equals $(e^{-\gamma}+o(1))/\log\log k$ as $k\to\infty$.
\end{corollary}

\begin{proof}
Replacing the distinct prime factors of $E$ by smaller primes, or adding factors up to $d(k)$ of them, can
only decrease the product in \cref{rec:thm:cyclic}; this gives the lower bound. For attainment put
$E=2\prod_{j\le d(k)}p_j$, choose a prime $Q\equiv1\pmod E$ and a character $\chi$ of order $E$ modulo $Q$. The parity
box has $2^{d(k)}<E$ distinct exponents, so a common additive shift avoids zero modulo $E$; every resulting
character is nonprincipal modulo the prime $Q$, hence primitive of conductor $Q$, and padding uses a
common primitive character. Common twisting and padding preserve local agreement. Finally
$\prod_{p\le p_d}(1-p^{-1})\sim e^{-\gamma}/\log p_d$ by Mertens' product theorem \cite[Eq.~(3)]{KeliherLee2024Constants}, $p_d\sim d\log d$, and
$d(k)\sim\log k/\log2$, so $\log p_{d(k)}\sim\log\log k$.
\end{proof}

\paragraph{Abelian $\ell$-groups.}
Let $\ell$ be a prime and $X\cong\bigoplus_{i=1}^r\Z/\ell^{a_i}\Z$ with all $a_i\ge1$; put $b=\min_ia_i$, $N=\abs X$, $M=N/\ell^b$,
$X[\ell]=\set{x:\ell x=0}$, and write $V$ additively, so that $(\ell v)(x)=v(x)^\ell$. Write $S=\mathcal D(c)$ and $\delta=\abs S/N$.
For $f(v)=\sum_x\widehat f(x)v(x)$ put $\norm f_A=\sum_x\abs{\widehat f(x)}$.

\begin{lemma}[Corrected anchored projection; \stU]\label{rec:lem:pproj}
For $y\in X$ and $a\in X$ put
\[
 K_y(v)=\mathbf 1_{v(y)=1}-\tfrac1\ell\mathbf 1_{v(\ell y)=1},\qquad J_a(v)=\sum_xc(x)K_{x-a}(v).
\]
Then $\mathbb E_vJ_a=(1-1/\ell)c(a)$, $\abs{J_a}\le m\mathbf 1_S$, and $\norm{J_a}_A\le2(1-1/\ell)m$.
\end{lemma}

\begin{proof}
For $y\ne0$ of order $\ell^j$, orthogonality gives $\mathbb EK_y=\ell^{-j}-\ell^{-1}\ell^{-(j-1)}=0$ and
$K_y(v)=\ell^{-j}\sum_{0\le u<\ell^j,\ \ell\nmid u}v(uy)$; for $y=0$, $K_y=1-1/\ell$. This gives the mean, and the Fourier
coefficients of each $K_y$ are nonnegative with sum $1-1/\ell$, which gives the norm bound. The values of $K_y$
are $1-1/\ell$, $-1/\ell$ and $0$, a range of length one; the two signed parts of $c$ each have mass $m$, and $J_a$ is
a sum of fibre masses of the value histogram, so it vanishes off $S$.
\end{proof}

\begin{theorem}[Sharp mass bound and its equality cases; \stU]\label{rec:thm:pmass}
For distinct packets of equal size, $\delta\ge(\ell-1)\norm c_\infty/(\ell m)\ge(\ell-1)/(\ell m)$. Equality $\delta=(\ell-1)/(\ell m)$
holds if and only if there are a subgroup $P\le X$, $a\in X$ and $t\in X[\ell]\setminus P$ with
$c=\mathbf 1_{a+P}-\mathbf 1_{a+t+P}$, so $m=\abs P$. The translate must have a representative of order $\ell$ in $X$; order $\ell$ in
$X/P$ does not suffice.
\end{theorem}

\begin{proof}
Average \cref{rec:lem:pproj} at an anchor where $\abs{c(a)}$ is largest, exchanging the packets if necessary. At
equality every nonzero coefficient is $\pm1$; let $A_0,B_0$ be the positive and negative supports, each of
size $m$. At a positive anchor, equality of the averages in $J_a\le m\mathbf 1_S$ forces pointwise equality on $S$,
so for $v\in S$ all elements of $A_0$ have one value under $v$, and the ratio of the value of each element of $B_0$
to that value is a nontrivial $\ell$th root of unity; with a negative anchor, all elements of $B_0$ also have one
value, and the two values differ by a root of exact order $\ell$. Let $P$ be generated by $A_0-A_0$ and $B_0-B_0$.
Then $S\subset P^\perp$, and $H_c$ vanishes off $S$ and equals $m(v(a)-v(b))\ne0$ on $S$ ($a\in A_0$, $b\in B_0$), so its support
is exactly $S$. Fourier inversion makes $c$ invariant under $P$, and as $A_0\subset a+P$, $B_0\subset b+P$, we get
$A_0=a+P$, $B_0=b+P$.

We claim $b-a\in P+X[\ell]$. Otherwise duality gives $v\in(P+X[\ell])^\perp$ with $v(b-a)\ne1$; since $X[\ell]^\perp=\ell V$, write
$v=\ell w$. The two coset histograms disagree at $v$, so $v\in S$, and then $w\in S$, since agreement at $w$ would
persist after raising all values to the $\ell$th power. But the description of $S$ makes $w(b-a)$ an $\ell$th root of
unity, contradicting $v(b-a)\ne1$. So $b-a=h+t$ with $h\in P$ and $t\in X[\ell]\setminus P$.

Conversely, for the coset pair and any $v$, the image $v(P)$ is a cyclic $\ell$-group and the image multiset of
$P$ is uniform on it. If $v(P)\ne1$ it contains all $\ell$th roots of unity, so multiplication by $v(t)$ preserves
the multiset; if $v(P)=1$, disagreement means $v(t)\ne1$. Hence $S=P^\perp\setminus\langle P,t\rangle^\perp$ and $\abs S=(1-1/\ell)N/\abs P$.
\end{proof}

\begin{theorem}[Saturation and prime-power sizes; \stU]\label{rec:thm:psat}
Every distinct equal-size pair has $\delta\ge\varphi(\ell^b)/N=(\ell-1)/(\ell M)$, and this bound is attained at size $k$
exactly when $k\ge M$. For every $k=\ell^u$,
\[
 d_X(k)=\frac{\ell-1}{\ell\min(k,M)},
\]
and for $k=\ell^u\le M$ the minimizers are exactly the pairs of \cref{rec:thm:pmass} with $\abs P=k$.
\end{theorem}

\begin{proof}
The set $S$ is invariant under multiplication by integers prime to $\ell$ (Galois conjugacy) and satisfies
$\ell v\in S\Rightarrow v\in S$; it is nonempty, since otherwise all traces vanish and $c=0$. Starting in $S$, take $\ell$th
roots as long as possible; this ends at $v_0\in S\setminus\ell V$. Such $v_0$ has a coordinate prime to $\ell$ in some cyclic
factor, so its order is at least $\ell^b$, and its unit orbit, of size at least $\varphi(\ell^b)$, lies in $S$. At equality $S$
is exactly the set of generators of $\langle v_0\rangle$ with $v_0\notin\ell V$ of order $\ell^b$; the transform of $c$ is supported
there, so $c$ is invariant under $\ker v_0$, of size $M$, and its positive mass is at least $M$. For
attainment split $X=C_{\ell^b}\oplus X'$ with $\abs{X'}=M$, take $P=X'$ and $t$ of exact order $\ell$ in the first factor, and
pad to every $k\ge M$. For $k=\ell^u\le M$ take instead a subgroup $P\le X'$ of order $k$; the coset pair attains
$(\ell-1)/(\ell k)$, the bound of \cref{rec:thm:pmass}, which leaves no room for cancellation, and the classification
follows.
\end{proof}

\begin{proposition}[Saturated equality; \stU]\label{rec:prop:psateq}
Equality in the saturation bound holds exactly for the nonzero balanced integral $c$ of the following
form: for some $v_0\in V\setminus\ell V$ of order $\ell^b$, $c$ is constant on the fibres of $v_0:X\to\mu_{\ell^b}$, and the induced
function $g$ on $C_{\ell^b}$ satisfies $\sum_{j=0}^{\ell-1}g(u+j\ell^{b-1})=0$ for every $u$ modulo $\ell^{b-1}$.
\end{proposition}

\begin{proof}
Necessity: by the preceding proof the transform of $c$ is supported on the generators of $\langle v_0\rangle$, and Fourier
inversion on $C_{\ell^{b-1}}$ says that the nonprimitive coefficients of $g$ vanish exactly when the displayed fibre
sums do. Conversely, for such a nonzero integral $g$ the primitive coefficients cannot vanish, since by
Galois conjugacy they would all vanish and then every coefficient would. If some multiple $dv$ lies in
the generator orbit, primitivity of $v_0$ forces $\ell\nmid d$, so $v$ lies in the same orbit; all power traces vanish
elsewhere, hence so do all value histograms, and the disagreement set is the orbit.
\end{proof}

\begin{remark}[Prime exponent]\label{rec:rem:pexp}
If $X$ has exponent $\ell$, then $b=1$, $X[\ell]=X$, $M=\ell^{r-1}$, every $v(\ell y)$ equals one and, as $\sum c=0$, $J_a$ is the
anchored projection of \cref{rec:lem:oradon}. \Cref{rec:thm:pmass} then gives the lowest level of \cref{rec:thm:olevel},
and \cref{rec:thm:psat,rec:prop:psateq} give the prime-power and saturated cases of \cref{rec:thm:odigit}; for elementary $2$-groups \cref{rec:thm:psat}
gives $2^{-\min(u+1,r)}$ at $k=2^u$, in accordance with \cref{rec:thm:qmatch}. For $X=(C_{\ell^a})^r$ it reads
$d_X(\ell^u)=(\ell-1)/\ell^{1+\min(u,a(r-1))}$; for $X=C_4\times C_4$ the minima at sizes $1,2,4,8,\dots$ are $1/2,1/4,1/8,1/8,\dots$.
The cyclic minimum of \cref{rec:cor:cglobal} tends to zero only like $1/\log\log k$, whereas elementary $2$-groups
of large rank allow disagreement $1/(2k)$ at every power of two $k$: the size alone does not govern the
possible ambiguity.
\end{remark}

\begin{example}[The torsion-lifting condition is needed]\label{rec:ex:ctorsion}
In $X=C_{\ell^2}$ take $P=\ell X$ and $t$ a generator. Although $t+P$ has order $\ell$ in $X/P$, it has no representative
of order $\ell$, and the packets $P$ and $t+P$ disagree at every nonidentity evaluation: a unit evaluation takes
the two cosets to disjoint cosets of $\mu_\ell$ in $\mu_{\ell^2}$, and a nonzero evaluation divisible by $\ell$ makes each
packet constant with distinct constants. So $\delta=1-\ell^{-2}$, not $(\ell-1)/\ell^2$.
\end{example}

\paragraph{Prime windows.}
Fix a smooth $w\ge0$ supported in $(\tfrac12,1)$ with $\int w=1$, and put $M_Y(f)=\sum_{p\nmid Q}(\log p)w(p/Y)f(v_p)$ for
$f:V\to\C$. Under Dirichlet GRH, \cref{rec:lem:gram} and Fourier expansion over $X$ give a constant $C_w$, depending
only on $w$, with
\begin{equation}\label{rec:eq:cgram}
 \abs{M_Y(f)-Y\,\mathbb Ef}\le C_w\norm f_A\sqrt Y\log(2Q)\qquad(Y\ge2)
\end{equation}
for every finite group $X$ of characters modulo $Q$ and every $f:V\to\C$. This is the only prime-sum estimate of
the section; \cref{rec:thm:qmatchfinite,rec:thm:ofinite} use it with $X$ of exponent two and $\ell$.

\begin{remark}[The standard scale]\label{rec:rem:cscale}
For every finite abelian $X$, whatever its exponent, the anchored trace $\Re(\overline{v(a)}H_c(v))$ has mean $c(a)$, norm
at most $2m$, and vanishes where the Euler factors agree; with \eqref{rec:eq:cgram} at a positive anchor it gives
a distinguishing prime in every window $(Y/2,Y)$ at the scale $Y\gg_wm^2\log^2(2Q)$. The finite theorems add the
best possible density and its extremal configurations, and for cyclic $X$ the possibly smaller scale
$4^t\log^2(2Q)$ of \cref{rec:thm:cwindow}; \cref{rec:prop:pwindow} gives the window bounds with $\norm c_\infty$ for every
abelian $\ell$-group.
\end{remark}

For cyclic $X$ and one primary
coordinate of order $\ell^a$, choose $h$ of order $\ell$; then $\mathbf 1_{v(h)\ne1}=1-\ell^{-1}\sum_{j<\ell}v(jh)$ has norm $2(1-1/\ell)$, and since
different primary coordinates have independent Fourier supports, the cylinders $C_I$ of the proof of
\cref{rec:thm:cyclic} satisfy, for every primary exponent,
\begin{equation}\label{rec:eq:cnorm}
 \norm{\mathbf 1_{C_I}}_A=2^{\abs I}d_I.
\end{equation}

\begin{theorem}[Cyclic window certificate; \stG]\label{rec:thm:cwindow}
Let $X$ be cyclic of order $E$, $t=t(E,k)$ and $D=D(E,k)$. For distinct packets of size $k$ and $Y\ge2$,
\[
 M_Y(\mathbf 1_{\mathcal D(c)})\ge DY\Bigl(1-C_w2^t\frac{\log(2Q)}{\sqrt Y}\Bigr),
\]
and the parity-box pairs of \cref{rec:thm:cyclic} have $\abs{M_Y(\mathbf 1_{\mathcal D(c)})-DY}\le C_w2^tD\sqrt Y\log(2Q)$. Hence every interval
$(Y/2,Y)$ contains a distinguishing prime once $Y\ge C_w'4^t\log^2(2Q)$; as $2^t\le2k$ this is $O(k^2\log^2(2Q))$, and it stops
growing with $k$ once $t=s$.
\end{theorem}

\begin{proof}
The proof of \cref{rec:thm:cyclic} gives $C_I\subseteq\mathcal D(c)$ with $\abs I\le t$ and $d_I\ge D$. Positivity of the weights,
\eqref{rec:eq:cgram} and \eqref{rec:eq:cnorm} give $M_Y(\mathbf 1_{\mathcal D(c)})\ge d_IY(1-C_w2^{\abs I}\log(2Q)/\sqrt Y)$. If the bracket with $2^t$ is
nonnegative, replace $\abs I$ by $t$ and $d_I$ by $D$; otherwise the claim is trivial. For the parity box, $\mathcal D(c)$ is the
cylinder on the first $t$ primary coordinates, and \eqref{rec:eq:cnorm} gives the two-sided bound.
\end{proof}

\begin{proposition}[Prime windows for $\ell$-groups; \stG]\label{rec:prop:pwindow}
Let $X$ be an abelian $\ell$-group and $\alpha=1-1/\ell$. For distinct packets of equal size,
\[
 M_Y(\mathbf 1_{\mathcal D(c)})\ge\frac\alpha m\bigl(\norm c_\infty Y-2C_wm\sqrt Y\log(2Q)\bigr).
\]
For the equality pairs of \cref{rec:thm:pmass,rec:prop:psateq}, $\norm{\mathbf 1_{\mathcal D(c)}}_A=2\alpha$ and
$\abs{M_Y(\mathbf 1_{\mathcal D(c)})-\delta(c)Y}\le2C_w\alpha\sqrt Y\log(2Q)$.
\end{proposition}

\begin{proof}
At an anchor $a$ with $c(a)=\norm c_\infty$ (after exchanging the packets), $J_a/m\le\mathbf 1_{\mathcal D(c)}$ by \cref{rec:lem:pproj}; it
may take negative values, but the prime weights are nonnegative, and \eqref{rec:eq:cgram} with the mean and norm
of \cref{rec:lem:pproj} gives the bound. In both equality families $\mathcal D(c)=H\setminus K$ for subgroups $K\le H\le V$ with
$[H:K]=\ell$ (by the proofs of \cref{rec:thm:pmass,rec:thm:psat}), and the coefficients of $\mathbf 1_H-\mathbf 1_K$ are positive on $H^\perp$,
negative on $K^\perp\setminus H^\perp$ and zero elsewhere, with positive and negative masses both $\alpha$.
\end{proof}

\begin{corollary}[Correction threshold; \stG]\label{rec:cor:ccorrect}
Let $X$ be cyclic of order $E$. If a record of local polynomials differs from each of two size-$k$ packets on
sets of upper relative density below $D(E,k)/2$, the packets coincide, and this threshold is sharp; these two
assertions are unconditional. For fixed $0<\eta<1$, with the weights $(\log p)w(p/Y)/M_Y(1)$, distinct size-$k$
packets are at weighted distance at least $(1-\eta)D(E,k)$ once $Y\ge C_w''\eta^{-2}4^t\log^2(2Q)$, so fewer than $(1-\eta)D(E,k)/2$
weighted errors leave at most one consistent packet.
\end{corollary}

\begin{proof}
The disagreement set of two packets lies in the union of their error sets. For sharpness split the
distinguishing primes of an extremal pair, in increasing order, into two alternating subsequences of
relative density $D/2$ each, and take the record from one packet on the first and from the other on the
second. In a window normalize by $M_Y(1)=Y+O_w(\sqrt Y\log(2Q))$ and apply \cref{rec:thm:cwindow}. This is a uniqueness
statement, not an efficient decoder.
\end{proof}

\begin{remark}[Plateaus]\label{rec:rem:cplateau}
The cyclic minimum is constant between successive powers of two in $k$ until all distinct prime factors
of $E$ are used:
\begin{center}
\begin{tabular}{@{}lccccc@{}}\toprule
$E$ & $k=1$ & $k=2,3$ & $k=4,\dots,7$ & $k\ge8$ & first saturation\\\midrule
$30$ & $1/2$ & $1/3$ & $4/15$ & $4/15$ & $4$\\
$210$ & $1/2$ & $1/3$ & $4/15$ & $8/35$ & $8$\\
$45$ & $2/3$ & $8/15$ & $8/15$ & $8/15$ & $2$\\
$72$ & $1/2$ & $1/3$ & $1/3$ & $1/3$ & $2$\\\bottomrule
\end{tabular}
\end{center}
Orders $6$, $12$, $72$ and $2^{100}3^{100}$ have the same matching function.
\end{remark}

\begin{example}[Two products of four $L$-functions modulo $31$]\label{rec:ex:c31}
Let $\chi$ be the character of order $30$ modulo $31$ with $\chi(3)=e^{2\pi i/30}$ ($3$ is a primitive root). The parity box on the
steps $15,10,6$ gives the exponent packets $\set{0,16,21,25}$ and $\set{1,6,10,15}$; twisting by $\chi^2$ gives
\[
 A'=\set{\chi^2,\chi^{18},\chi^{23},\chi^{27}},\qquad B'=\set{\chi^3,\chi^8,\chi^{12},\chi^{17}},
\]
all primitive of conductor $31$. The two products of four Dirichlet $L$-functions have different local
factors exactly at the primes that are primitive roots modulo $31$, of relative density $\varphi(30)/30=4/15$, so
they agree at $11/15$ of the primes, the largest agreement among distinct degree-four packets in this group
by \cref{rec:thm:cyclic}; the products themselves differ.
\end{example}

\begin{example}[Quartic characters modulo $65$]\label{rec:ex:c65}
Let $\alpha$ and $\beta$ have order four modulo $5$ and $13$, viewed modulo $65$, so $X=C_4\times C_4$. The packets
$\set{1,\alpha,\alpha^2,\alpha^3}$ and $\set{\beta^2,\beta^2\alpha,\beta^2\alpha^2,\beta^2\alpha^3}$ have different Euler factors exactly when $\alpha(p)=1$ and $\beta^2(p)=-1$, that is
for the six classes $6,11,21,31,41,46$ among the $48$ units modulo $65$; the disagreement density is $1/8$, optimal
in this group for every size at least four by \cref{rec:thm:psat}. At $p=11$ the two polynomials are $(1-T)^4$
and $(1+T)^4$.
\end{example}

Exact enumerations by integer histogram comparison cover $10{,}417{,}835$ unordered packet pairs
($2{,}453{,}747$ in $53$ cyclic cases, $7{,}964{,}088$ in nine $\ell$-group cases, among them all $7{,}509{,}750$ pairs for $C_4^2$ at
$k=4$, with minimum $1/8$ attained by exactly $12$ pairs), the equality predicates as equivalences, $49$ cyclic
constructions, $38$ cylinder transforms, $24$ subgroup constructions and both examples:
\begin{center}
\artifact{euler_agreement_2026-11-27/data/verification_summary.json}\\
\artifact{euler_agreement_2026-11-27/data/cyclic_euler_verification.json}\\
\artifact{euler_agreement_2026-11-27/data/pprimary_verification.json}
\end{center}
Related uncertainty principles and agreement results for finite-group characters are in
\cite{Meshulam2006Uncertainty,MartinWalji2017Distinguishing}.

\subsection{Matching Euler factors at odd prime order}\label{rec:sec:omatch}
\sources{euler-packet-rigidity; odd-order-euler-factor-matching; odd_prime_order_euler_packets; cubic_quintic_euler_agreement; odd_euler_agreement; odd_prime_euler_agreement; odd_prime_euler_agreement_2; odd_prime_euler_packets; prime_exponent_euler_agreement; prime_order_euler_agreement; grh_cubic_matching_and_precision_2026-11-27; gram-euler-and-zero-filters}

Let $\ell$ be an odd prime and $X$ a group of Dirichlet characters modulo $Q$ of exponent $\ell$ and rank
$r\ge1$ over $\F_\ell$, written additively, $V=X^*$, $\zeta_\ell=e^{2\pi i/\ell}$, and for $p\nmid Q$ let $v_p\in V$ be given by
$\chi(p)=\zeta_\ell^{v_p(\chi)}$. Packets $A,C$ are multisets of $k$ members of $X$, compared through
$\prod_{\chi\in A}(1-\chi(p)T)$, and $\delta(A,C)$ is the density of primes at which these local polynomials differ.
Put $c_x=\operatorname{mult}_A(x)-\operatorname{mult}_C(x)$, let $m=\sum_xc_x^+\le k$ be the mass of each side after cancellation,
\[
 H(v)=\sum_{x\in X}c_x\zeta_\ell^{v(x)},\qquad S=\set{v\in V:H(v)\ne0},\qquad \beta_m=\frac{\ell-1}{\ell m},
\]
and let $d_\ell(r,k)$ be the least $\delta(A,C)$ over distinct packets of size $k$. Common entries are allowed, and
a packet may lie in a proper subspace. The results below are unconditional except
\cref{rec:thm:ofinite}. The least disagreement density is the universal staircase bound of
\cref{rec:thm:odigit} exactly when the leading base-$\ell$ digit of $k$ divides $\ell-1$ or $\ell+1$ or the size is
saturated, so for every rank and size when $\ell=3,5$; in every such case all extremal pairs are
classified (\cref{rec:thm:oplanar,rec:thm:odigit}, and \cref{rec:thm:psat,rec:prop:psateq} for saturated sizes). For the other digits \cref{rec:thm:odescent} reduces the minimum
exactly to packets of size below $\ell$ in rank at least two, \cref{rec:thm:oconic} settles a range of them, and
the minimum is known for every rank and size when $\ell\le11$ (\cref{rec:cor:oseveneleven}). The results on
cosets, prime-power sizes and saturation hold for every finite abelian $\ell$-group (\cref{rec:rem:pexp}).

\begin{remark}[Traces decide]\label{rec:lem:otrace}
The exponent $\ell$ has $1$ as its only positive proper divisor, so by \cref{rec:prop:cdivisor} the local polynomials
agree at $p$ exactly when $H(v_p)=0$, and by \cref{rec:lem:chist}(iii) $\delta(A,C)=\abs S/\ell^r$, which depends only on
$(\ell,r,c)$. The set $S$ omits $0$ since $\sum_xc_x=0$, is a union of punctured lines through the origin by
\cref{rec:lem:chist}(ii), and is nonempty when $A\ne C$ by Fourier inversion. At exponent two the same remark,
with $E=2$, is the trace criterion of \cref{rec:sec:qmatch}.
\end{remark}

Every rank occurs: choose $r$ primes $q\equiv1\pmod\ell$ and an order-$\ell$ character at each. In rank one every
distinct pair has $\delta=1-1/\ell$.

\begin{remark}[Anchored projection count]\label{rec:lem:oradon}
At exponent $\ell$ the function $K_y$ of \cref{rec:lem:pproj} is $\mathbf 1_{v(y)=0}-1/\ell$, and as $\sum_xc_x=0$ its $J_a$ is the
signed mass of the projection fibre through $a$, $J_a(v)=\sum_xc_x\mathbf 1_{v(x-a)=0}$. So
\[
 \mathbb E_vJ_a=(1-1/\ell)c_a,\qquad\abs{J_a}\le m\mathbf 1_S,\qquad\norm{J_a/m}_A\le2(1-1/\ell),
\]
where $\norm f_A=\sum_{x\in X}\abs{\widehat f(x)}$ for $f(v)=\sum_x\widehat f(x)\zeta_\ell^{v(x)}$, and \cref{rec:thm:pmass} gives
$\delta(A,C)\ge\beta_m\norm c_\infty$. This fibre form is used in \cref{rec:thm:olevel,rec:thm:ofinite}.
\end{remark}

\begin{lemma}[Unions of hyperplanes; \stU]\label{rec:lem:ounion}
Let $r\ge2$ and let $H_1,\dots,H_d$ be distinct linear hyperplanes of $V$, $1\le d\le\ell$. Then
$\abs{\bigcup_iH_i}\le\ell^{r-1}+(d-1)(\ell^{r-1}-\ell^{r-2})$, with equality for $d\ge2$ exactly when all $H_i$ contain a
common subspace of codimension two. A linear hyperplane is never covered by $d\le\ell$ other linear
hyperplanes.
\end{lemma}

\begin{proof}
Add each $H_i\setminus H_1$, of size $\ell^{r-1}-\ell^{r-2}$. Equality means these pieces are disjoint, so
$H_i\cap H_j\subseteq H_1$ for $i,j>1$, and by dimensions $H_i\cap H_1=H_j\cap H_1$; the converse is clear. The traces of $d$
other hyperplanes on a hyperplane are proper subspaces containing $0$, so their union has at most
$1+d(\ell^{r-2}-1)<\ell^{r-1}$ points.
\end{proof}

Let $\OO=\Z[\zeta_\ell]$ and $\pi=\zeta_\ell-1$, so $\OO/(\pi)=\F_\ell$ and $(\ell)=(\pi)^{\ell-1}$. All $H(v)$ lie in $\pi\OO$; let $t$ be the
least $\pi$-adic valuation of a nonzero value. Write $a=\floor{\log_\ell k}$ and $b=\floor{k/\ell^a}\in\set{1,\dots,\ell-1}$.

\begin{lemma}[The first nonzero reduction; \stU]\label{rec:lem:ored}
If $t=j(\ell-1)+s$ with $j\ge0$ and $1\le s\le\ell-1$, then $k\ge\ell^js$; hence $t\le a(\ell-1)+b$. The function
$f(v)=\pi^{-t}H(v)\bmod\pi$ is a nonzero polynomial function on $V$ of degree at most $t$ with $f(0)=0$ and
$f(uv)=u^tf(v)$ for $u\in\F_\ell^\times$.
\end{lemma}

\begin{proof}
Evaluation at $\zeta_\ell$ is an additive isomorphism from the zero-sum vectors of $\Z^\ell$ onto $\pi\OO$ (injective by
the cyclotomic argument above, surjective since the $\zeta_\ell^{i+1}-\zeta_\ell^i$ form a $\Z$-basis). At $v$ attaining $t$,
$H(v)\in\ell^j\pi\OO$, so the histogram difference is $\ell^jn$ for an integral zero-sum $n$ whose positive and
negative parts have mass $m'\le k/\ell^j$. Expanding $\zeta_\ell^i=(1+\pi)^i$, valuation $s$ forces
$\sum_in_i\binom i\lambda\equiv0\pmod\ell$ for $\lambda<s$, so the two $m'$-element multisets represented by $n^\pm$ have equal power
sums through degree $s-1$; if $m'<s$, Newton's identities (valid as $m'<\ell$) make them equal, a
contradiction. So $k\ge\ell^jm'\ge\ell^js$, and the numbers $\ell^js$ increase with $j(\ell-1)+s$. For $f$, lift coordinates
to $\set{0,\dots,\ell-1}$ and expand $\zeta_\ell^{v(x)}=\prod_i\sum_{u_i}\binom{v_i}{u_i}(\zeta_\ell^{x_i}-1)^{u_i}$: the coefficient of the tensor
binomial indexed by $u$ is divisible by $\pi^{\abs u}$, and, being an iterated forward difference of values of $H$,
by $\pi^t$. Dividing by $\pi^t$ and reducing kills the terms with $\abs u>t$. Finally $\sigma_u(\pi)/\pi\equiv u\pmod\pi$ for
the automorphism $\sigma_u:\zeta_\ell\mapsto\zeta_\ell^u$, which acts trivially on the residue field, so $f(uv)=u^tf(v)$.
\end{proof}

For $n,D\ge0$ put $W(n,D)=\min\set{\prod_{i=1}^n(\ell-e_i):0\le e_i\le\ell-1,\ \sum e_i\le D}$. Moving a unit between unsaturated
exponents shows $W(n,D)=1$ if $D\ge n(\ell-1)$, and otherwise $W(n,D)=(\ell-h)\ell^{n-j-1}$ for $D=j(\ell-1)+h$, $0\le h<\ell-1$;
$W(n,\cdot)$ decreases strictly until it reaches $1$.

\begin{lemma}[Surviving directions; \stU]\label{rec:lem:oproj}
A nonzero polynomial function of degree at most $D$ in $n$ variables over $\F_\ell$ has at least $W(n,D)$
nonzero values. If a nonzero polynomial function on $\F_\ell^{n+1}$ of degree at most $t\ge1$ vanishes at $0$
and its support is a union of punctured lines, its support contains at least $W(n,t-1)$ projective
directions.
\end{lemma}

\begin{proof}
Reduce exponents below $\ell$. If $e_1$ is the top exponent of the first variable, above each point where its
coefficient is nonzero at least $\ell-e_1$ values survive; iterate on that coefficient. For the second claim
induct on $n$, the case $n=0$ being clear. If some projective hyperplane, say $x_0=0$, misses the support,
then $F(0,x')$ vanishes identically and $F=x_0G$ with $\deg G\le t-1$; the chart $x_0=1$ meets every supported
line once, and the first claim applies. Otherwise induction in each of the $(\ell^{n+1}-1)/(\ell-1)$ hyperplanes and
double counting give more than $\ell W(n-1,t-1)\ge W(n,t-1)$ directions.
\end{proof}

\begin{construction}[Extremal packets]\label{rec:con:opack}
In rank two:
\emph{split}, if $b\mid\ell-1$: for $H_0\le\F_\ell^\times$ of order $b$, the packets $H_0e_1$ and $H_0e_2$ agree at $(u,w)\ne0$
exactly when $uw\ne0$ and $u/w\in H_0$;
\emph{nonsplit}, if $b\mid\ell+1$: identify the plane with $K=\F_{\ell^2}$ and its dual through $\Tr_{K/\F_\ell}(vx)$, let $U$ be the
norm-one group, $H_0\le U$ of order $b$ and $c_0\in U\setminus H_0$; for $v\ne0$ all elements of $vH_0$ and $vc_0H_0$ have the same norm,
equal traces and norms force equality or Frobenius conjugacy, so the packets $H_0$, $c_0H_0$ agree exactly
when $v^{\ell-1}\in c_0H_0$, and $v\mapsto v^{\ell-1}$ is a bijection $K^\times/\F_\ell^\times\to U$.
Either way exactly $b$ of the $\ell+1$ directions agree. \emph{Cosets}: for a subspace $U'\le X$ of dimension $s$ and
$y-x\notin U'$, $H(v)=\ell^s\mathbf 1_{U'^\perp}(v)(\zeta_\ell^{v(x)}-\zeta_\ell^{v(y)})$ and $\delta(x+U',y+U')=(\ell-1)/\ell^{s+1}$. \emph{Lifting}: replacing every
entry $x$ by $x+L$ for a subspace $L$ of dimension $a$ independent of the support multiplies $H$ by
$\ell^a\mathbf 1_{L^\perp}$ and divides $\delta$ by $\ell^a$. Embedding a plane in $X$, padding both sides with common entries and a
common translation preserve agreement; in a plane a translation outside the at most $2k<\ell^2$ points of
$-(A\cup C)$ makes every factor of exact order $\ell$.
\end{construction}

In an affine plane $E\cong\F_\ell^2$ of $X$, a \emph{matching direction} of two finite multisets is a
one-dimensional space of linear forms on $E$ under which their projection multisets coincide; for
packets in $E$ these are exactly the projective directions of agreement.

\begin{theorem}[Planar switching configurations; \stU]\label{rec:thm:oplanar}
Let $1\le b\le\ell-1$ and let $A,C\subset\F_\ell^2$ be disjoint sets of $b$ points. They have at most $b$ matching
directions, and $b$ is attained if and only if $b\mid\ell-1$ or $b\mid\ell+1$. For $b=1$ any two points attain
it; for $b=2$ exactly the opposite vertex pairs $\set{x,x+u+w}$, $\set{x+u,x+w}$ of a parallelogram ($u,w$
independent); for $b\ge3$ exactly the following configurations, up to an affine change of coordinates,
where $G$ is a subgroup of order $b$ of the indicated cyclic group:
\begin{enumerate}
\item $A=G\times\set0$ and $C=\set0\times G$ with $G\le\F_\ell^\times$ (two intersecting lines, their common point omitted);
\item $A=\set{(t,t^{-1}):t\in G}$ and $C=\set{(t,t^{-1}):t\in\lambda G}$ with $G\le\F_\ell^\times$ and $\lambda\notin G$ (a split hyperbola);
\item $A=G$ and $C=\lambda G$ in $\F_\ell^2\cong\F_{\ell^2}$, with $G$ in the norm-one group of $\F_{\ell^2}$ and $\lambda\notin G$ of norm one
 (a nonsplit conic).
\end{enumerate}
\end{theorem}

\begin{proof}
\emph{Counting.} Fix $a\in A$. At a matching direction with kernel line $\kappa$ some $c\in C$ has $c-a\in\kappa$, and
distinct matching directions have distinct kernels, so they use distinct points of $C$; there are at
most $b$. At equality every cross difference $c-a$ lies on a matching kernel, from each point of either
set each matching kernel contains exactly one cross difference, and so every matching projection is
injective on $A$ and on $C$ (a repeated value on $A$ would put two points of $C$ on one kernel line through
$a$). For $b=1$ nothing is to prove. For $b=2$ take two matching forms as coordinates: $A=\set{(x_1,y_1),(x_2,y_2)}$
with $x_1\ne x_2$, $y_1\ne y_2$, and $C$, disjoint from $A$ with the same coordinate values, is $\set{(x_1,y_2),(x_2,y_1)}$; this
is a parallelogram, and $2\mid\ell-1$.

\emph{Moments.} Let $b\ge3$ and $R_A(Z;u,v)=\prod_{(x,y)\in A}(Z-ux-vy)$, similarly $R_C$. The difference of the
coefficients of $Z^{b-j}$ is a binary form of degree $j$ vanishing at the $b$ matching directions, hence zero
for $j<b$. By Newton's identities the power sums of $ux+vy$ over $A$ and $C$ agree through degree $b-1$, and
since binomial coefficients below $\ell$ are invertible,
\begin{equation}\label{rec:eq:omixed}
 \sum_{(x,y)\in A}x^iy^j=\sum_{(x,y)\in C}x^iy^j\quad\text{in }\F_\ell\qquad(i+j\le b-1).
\end{equation}

\emph{The conic.} Take two matching forms as coordinates $x,y$. By injectivity, $A$ and $C$ are graphs
$\set{(s,f(s))}$, $\set{(s,g(s))}$ over one $b$-set $S\subset\F_\ell$, with $f(s)\ne g(s)$. Put $P(X)=\prod_{s\in S}(X-s)$. Since
$\sum_{s\in S}F(s)/P'(s)$ is the coefficient of $X^{b-1}$ in the interpolation polynomial of $F$ on $S$, a function
$h$ on $S$ with $\sum_ss^ih(s)/P'(s)=0$ for $0\le i\le b-m-2$ has interpolation polynomial of degree at most $m$
(if the degree were $d>m$, take $i=b-1-d$). By \eqref{rec:eq:omixed} with $y$-degree $1,2,3$ this gives
$f-g=\gamma/P'$ with $\gamma\in\F_\ell^\times$, then $f+g=L$ with $\deg L\le1$, and $f^2+fg+g^2=R$ with $\deg R\le2$ (for $b=3$ by
interpolation at three points). With $Y=y-L(x)/2$ both sets lie on $Y^2=D(x)$, $D=R-\tfrac34L^2$, $\deg D\le2$, and
above each $s\in S$ the two sets take the opposite values $\pm(f-g)(s)/2$; so $D(s)$ is a nonzero square.

\emph{The cases.} If $S_1,S_2$ are subsets of an abelian group with $\abs{S_1}=\abs{S_2}=\abs{S_1S_2}=b$, then all translates
$S_1s$, $s\in S_2$, coincide, so the stabilizer of $S_1$ has order $b$ and both sets are cosets of it. If $\deg D=1$
the conic is a parabola $(t,t^2)$, the cross direction of parameters $u\ne v$ is that of $(1,u+v)$, and the
additive group of order $\ell$ would need a subgroup of order $b$ with $1<b<\ell$. If $D$ has two distinct roots,
completing the square gives $Y^2-\alpha x'^2=\beta$ with $\alpha\beta\ne0$. For square $\alpha$ this is the hyperbola $(t,t^{-1})$, the
cross direction is that of $(1,-1/(uv))$, determined injectively by $uv$, and the parameter sets are cosets of
an order-$b$ subgroup of $\F_\ell^\times$: case (ii), $b\mid\ell-1$. For nonsquare $\alpha$ the conic is a norm level of
$\F_{\ell^2}$, scaled to the norm-one group; for norm-one $u\ne v$, $(u-v)^{\ell-1}=-1/(uv)$, and $z\mapsto z^{\ell-1}$ is injective on
directions, so we get case (iii), $b\mid\ell+1$. Otherwise $D$ is a nonzero square constant or a square times
$(x-x_0)^2$ (a nonsquare multiple would give nonsquare or zero values on $S$), so the conic is two parallel
or two intersecting lines. If both lines carried points of $A$, an $A$-point on the first line would have
as many cross differences along that line as there are $A$-points on the second; injectivity forces both
numbers to be at most one, against $b\ge3$. So $A$ and $C$ lie on different lines. For parallel lines the
cross directions are determined by differences of parameters, excluded as before. For intersecting lines
the common point lies in neither set (from it all cross differences would be parallel); on the axes the
cross directions are determined by the ratios of the parameters, and rescaling the axes gives case (i),
$b\mid\ell-1$.

\emph{Sufficiency.} In each of (i)--(iii) the cross direction is an injective function of the ratio,
respectively the product, of the parameters, which takes $b$ values, and from each point each value is
taken once; the induced bijection $A\to C$ preserves the corresponding projection, so these $b$ directions
match, and no others by the count. Disjoint cosets exist in (ii), (iii) because $b$ is smaller than the
group order there; for $b=\ell-1$ use (i).
\end{proof}

The configurations with $b\ge3$ are the dual $3$-nets with a collinear component, in accordance with
\cite[Theorem~5.1 and \S5]{BlokhuisKorchmarosMazzocca2011Structure}; the proof above needs only interpolation.

\begin{theorem}[Packets below the character order; \stU]\label{rec:thm:osmall}
Let $r\ge2$ and $1\le k\le\ell$, and put $D_1(k)=(\ell-1)(\ell+1-k)/\ell^2$.
\begin{enumerate}
\item Every distinct pair has $\delta\ge D_1(k)$.
\item For $2\le k\le\ell$ a pair attains equality if and only if nothing cancels and $A,C$ are disjoint $k$-sets in
a common affine plane of $X$ whose cross differences $x-y$, $x\in A$, $y\in C$, occupy exactly $k$ projective
directions, each exactly once from every point of either set. The agreement set is then a pencil of $k$
hyperplanes, and $A$, $C$ and the $k$ directions at infinity form an embedded dual $3$-net of order $k$.
\item Equality is attained, in every rank $r\ge2$, if and only if $k=1$, $k=\ell$, $k\mid\ell-1$ or $k\mid\ell+1$. For
$2\le k\le\ell-1$ the equality pairs are, in their plane, exactly the configurations of \cref{rec:thm:oplanar}; for
$k=\ell$ they are two parallel affine lines.
\item If $3\le k<\ell$ and $k$ divides neither $\ell-1$ nor $\ell+1$, every distinct pair of size $k$ has
$\delta\ge(\ell-1)(\ell+2-k)/\ell^2$, with equality attained when $k-1$ divides $\ell-1$ or $\ell+1$.
\end{enumerate}
\end{theorem}

The criterion in (iii) is the same in every rank; for $k=2,3,4$ it always holds, and $k=\ell-1$ gives
$2(\ell-1)/\ell^2$.

\begin{lemma}[Separating a point; \stU]\label{rec:lem:osep}
Let $n\ge2$ and let $S_0$ be a set of at most $2n+2$ points of an affine space over a field, contained in no affine
plane. There are $P\in S_0$ and a polynomial of degree at most $n$ vanishing on $S_0\setminus\set P$ and not at $P$.
\end{lemma}

\begin{proof}
If a line contains at least $n+2$ points of $S_0$, take $P$ off it, one hyperplane through the line avoiding $P$
and one hyperplane avoiding $P$ through each of the at most $n-1$ remaining points. Otherwise fix $P$ and
group the vectors $x-P$ by their spans; each group has at most $n$ members and at most two have $n$. Choose
an independent triple meeting every group of size $n$ (the vectors have rank at least three); the at most
$2(n-1)$ others, in groups of size at most $n-1$, split into at most $n-1$ independent pairs or singletons
(pad to $2(n-1)$ entries in new groups, sort by group and pair entries $n-1$ places apart). For each part
choose a linear form equal to $1$ on its vectors; the product of the at most $n$ affine equations works.
\end{proof}

\begin{proof}[Proof of \cref{rec:thm:osmall}]
(i) Cancel, and fix $x\in A$. Agreement at $v$ forces $v(x)=v(y)$ for some $y\in C$, so the agreement set lies in
at most $d\le m\le k$ distinct hyperplanes $(x-y)^\perp$, and \cref{rec:lem:ounion} gives $\delta\ge(\ell-1)(\ell+1-d)/\ell^2$.
(ii) At equality $d=m=k$, so nothing cancels, $C$ is a set, the union is the whole agreement set and is a
pencil with axis $W'$ of codimension two; by symmetry $A$ is a set. By the last clause of
\cref{rec:lem:ounion} the $k$ hyperplanes of every anchor are members of this pencil, so all cross differences
lie in the plane $W'^\perp\subset X$ with the $k$ normal directions, each once from each anchor, and fixing one point of
each set puts both in one translate of $W'^\perp$. Conversely, each cross direction then matches $A$ with $C$, so
the $k$ hyperplanes it defines all agree, and no other $v$ agrees by the inclusion; the pencil attains the
bound. With the directions at infinity, every line meeting two of the three sets meets each once.
(iii) Singletons, the coset construction with $s=1$, and the split and nonsplit constructions give
sufficiency. By (ii) an equality pair with $2\le k\le\ell-1$ consists of disjoint $k$-sets in an affine plane $E$ of $X$
with exactly $k$ matching directions (the restriction of $v$ to $E$ is uniform), so \cref{rec:thm:oplanar} gives
necessity and the configurations. For $k=\ell\le\ell^{r-1}$ the minimum is $(\ell-1)/\ell^2$ and \cref{rec:thm:psat} leaves only the
coset pairs of \cref{rec:thm:pmass} with $\abs P=\ell$, two cosets of one line.
(iv) If the support $S_0$ of $c$ (at most $2k$ points) lies in an affine plane, translate it into a linear plane;
restriction of $v$ to that plane is uniform, the disagreement count there is a multiple of $\ell-1$, and strict
inequality in (i) gives the bound. Otherwise \cref{rec:lem:osep} with $n=k-1$ gives a polynomial $F$ of degree at
most $k-1$ with $\sum_xc_xF(x)\ne0$ in $\F_\ell$, hence, as $\sum c_x=0$, a monomial $x^\alpha$ of degree $1\le j\le k-1$ with
$\sum_xc_xx^\alpha\ne0$. Then $M_j(v)=\sum_xc_xv(x)^j$ has coefficient $\binom j\alpha\sum_xc_xx^\alpha\ne0$ at $v^\alpha$, since $j<\ell$; it
vanishes wherever the projections agree, and \cref{rec:lem:oproj} gives at least $W(r-1,j-1)$ disagreeing
directions, density at least $(\ell-1)(\ell+1-j)/\ell^2\ge(\ell-1)(\ell+2-k)/\ell^2$. Pad an extremal pair of size $k-1$ for the
equality.
\end{proof}

\begin{example}[Order seven]\label{rec:cor:oseven}
For $\ell=7$ and $k=5$, five divides neither $6$ nor $8$ and four divides $8$, so \cref{rec:thm:osmall}(iv) gives $d_7(r,5)=24/49$ for every
$r\ge2$: two distinct products of five septic $L$-functions agree at no more than $25/49$ of the primes, and
this is attained. In rank two the positive disagreement densities of size-$5$ pairs are exactly
$\set{24,30,36,42,48}/49$. Indeed $\delta=(\ell-1)(\ell+1-d)/\ell^2$ with $d\le m\le k$ the number of agreeing directions;
$d=5$ is excluded by \cref{rec:thm:osmall}(iii), $d=1,\dots,4$ are realized at size $d$ and padded, and $d=0$ by two copies of $0$
against two independent points, padded.
\end{example}

\begin{theorem}[The leading base-$\ell$ digit; \stU]\label{rec:thm:odigit}
For distinct packets of size $k$,
\[
 \delta(A,C)\ge D_\ell(r,k)=\begin{cases}(\ell-1)(\ell+1-b)/\ell^{a+2},&a\le r-2,\\ (\ell-1)/\ell^r,&a\ge r-1.\end{cases}
\]
\begin{enumerate}
\item If $a\ge r-1$ the bound is the saturation bound of \cref{rec:thm:psat} \textup($b=1$, $M=\ell^{r-1}$\textup), the exact minimum,
and its equality cases are those of \cref{rec:prop:psateq}.
\item If $a\le r-2$ it is the exact minimum if and only if $b\mid\ell-1$ or $b\mid\ell+1$. In every equality case there is a
subspace $L\le X$ of dimension $a$ such that, after cancellation, both packets are full inverse images under
$X\to X/L$ of disjoint $b$-sets in a common affine plane with $b$ matching directions, that is, two points for
$b=1$, the opposite vertex pairs of a parallelogram for $b=2$, and for $b\ge3$ one of the three configurations
of \cref{rec:thm:oplanar}; in particular the reduced size is exactly $b\ell^a$ and the reduced packets are sets.
Conversely every such lift, padded to size $k$, attains $D_\ell(r,k)$.
\item For $k=\ell^a$ with $a\le r-1$ the minimum and the minimizers are given by \cref{rec:thm:psat,rec:thm:pmass}, as
$X[\ell]=X$. Equality in $\delta\ge(\ell-1)/(\ell k)$ forces $k$ to be a power of $\ell$, by \cref{rec:thm:pmass} with $m\le k$.
\item $D_\ell(r,k)=d_\ell(r,k)$ for every $r$ and $k$ when $\ell=3$ or $\ell=5$.
\end{enumerate}
\end{theorem}

\begin{proof}
By \cref{rec:lem:ored,rec:lem:oproj} with $n=r-1$, $S$ contains at least $W(r-1,t-1)\ge W(r-1,a(\ell-1)+b-1)$ directions, each
with $\ell-1$ points, which is the bound. (ii) Sufficiency: lift the rank-two construction for $b$ by an $a$-dimensional $L$ and pad. Necessity: as
$W(r-1,\cdot)$ is strictly decreasing here, equality forces $t=e:=a(\ell-1)+b$ and $S=\supp f$. By \cref{rec:lem:ored}
the reduced representative of $f$ has only monomials of degree $\equiv e\pmod{\ell-1}$ and no constant term;
raising one positive exponent of each makes it homogeneous of degree $e$ with the same values. So $f$ is a
minimum-weight word of the projective Reed--Muller code of degree $e$ on $\mathbf P^{r-1}(\F_\ell)$, and by
\cite[Thm.~4.3]{GhorpadeLudhani2024Minimum} it has a factor $L_a\prod_{i<a}(L_a^{\ell-1}-L_i^{\ell-1})$ with independent linear
forms; hence $S$ lies in the common kernel $U$ of $L_0,\dots,L_{a-1}$, of codimension $a$. Fourier inversion makes $c$
invariant under $L=U^\perp$; the induced $c'$ on $X/L$ has masses at most $\floor{k/\ell^a}=b$ and density $\ell^aD_\ell(r,k)=D_1(b)$,
so \cref{rec:thm:osmall}(ii) applies: $c'$ is the difference of the indicators of two disjoint $b$-sets in a plane
with $b$ matching directions, classified by \cref{rec:thm:oplanar}. (iv) The admissible digits $1,2$ for $\ell=3$ and $1,2,3,4$ for $\ell=5$ all divide $\ell\pm1$.
\end{proof}

For example, the least disagreement densities at rank two, for $k=1,\dots,7$ (constant from $k=\ell$ on), are
$2/3,4/9,2/9,\dots$ for $\ell=3$; $4/5,16/25,12/25,8/25,4/25,\dots$ for $\ell=5$; and
$6/7,36/49,30/49,24/49,24/49,12/49,6/49$ for $\ell=7$.

\begin{remark}[Cubic packets]\label{rec:cor:ocubicmin}
For $\ell=3$ every leading digit divides $\ell\pm1$ (\cref{rec:thm:odigit}(iv)), so with $a=\floor{\log_3k}$ the least
disagreement density of two distinct size-$k$ packets of rank $r$ is
\[
 \Delta_{r,k}=\begin{cases}2/3^{a+1},&a\le r-2,\ 3^a\le k<2\cdot3^a,\\ 4/3^{a+2},&a\le r-2,\ 2\cdot3^a\le k<3^{a+1},\\ 2/3^r,&k\ge3^{r-1}.\end{cases}
\]
Up to exchanging the packets and adding a common multiset, the extremal pairs of \cref{rec:thm:odigit}(ii) and
\cref{rec:prop:psateq} take the following normal form. In the first range they are $\eta+U$ and $\eta+\theta+U$ with
$\dim U=a$, $\theta\notin U$. In the second range the lift of a parallelogram $\set{x,x+u+w}$, $\set{x+u,x+w}$ becomes, with
$\eta=x+(u+w)/2$, $\alpha=(u+w)/2$, $\beta=(u-w)/2$ ($2$ is invertible), the pair $(\eta+\alpha+U)\sqcup(\eta-\alpha+U)$ and
$(\eta+\beta+U)\sqcup(\eta-\beta+U)$ with $\dim U=a$ and $\alpha,\beta$ independent modulo $U$. In the saturated range $c$ is
constant on the three cosets of a hyperplane, with values $d_0+d_1+d_2=0$, not all zero, and $3^{r-1}\max_j\abs{d_j}\le k$,
since a zero-sum triple has positive mass $\max_j\abs{d_j}$. The pairs $\pm\theta+U$, the pairs with $\eta=0$ and the two
nonzero cosets of a hyperplane avoid the principal character, and padding can use a nonprincipal one, so $\Delta_{r,k}$
is also the minimum when every entry is required to be nonprincipal. As in the proof of \cref{rec:cor:ccorrect},
products of $k$ characters of $X$ observed through their Euler factors at the primes outside a set of upper
density below $\Delta_{r,k}$ are determined, and a record of local factors that differs from the true one on a set of
upper density below $\Delta_{r,k}/2$ has at most one packet within that distance; both thresholds are sharp.
\end{remark}

For example, in rank three two products of six cubic $L$-functions agree at most at $23/27$ of the primes,
nine factors reach $25/27$, and in rank at least five two products of $54$ factors can agree at $239/243$ of the
primes, in which case each side has, after cancellation, exactly $54$ distinct characters.

\begin{example}[Six cubic factors modulo $1729$]\label{rec:ex:ocubic}
Let $Q=7\cdot13\cdot19=1729$ and let $\chi_7,\chi_{13},\chi_{19}$ be the cubic characters with $\chi_7(3)=\chi_{13}(2)=\chi_{19}(2)=\zeta_3$.
The products
\[
 \prod_{j=0}^2L(s,\chi_7^j\chi_{13})L(s,\chi_7^j\chi_{13}^{-1})\quad\text{and}\quad\prod_{j=0}^2L(s,\chi_7^j\chi_{19})L(s,\chi_7^j\chi_{19}^{-1})
\]
have six distinct nonprincipal factors each, and
\[
 H(p)=\bigl(1+\chi_7(p)+\chi_7(p)^2\bigr)\bigl(\chi_{13}(p)+\overline{\chi_{13}(p)}-\chi_{19}(p)-\overline{\chi_{19}(p)}\bigr)
\]
is nonzero exactly when $\chi_7(p)=1$ and
exactly one of $\chi_{13}(p),\chi_{19}(p)$ is $1$. So the Euler factors agree at $23/27$ of the primes, the maximum by
\cref{rec:cor:ocubicmin}. Exactly $192$ of the $1296$ reduced classes modulo $1729$ disagree; the first
distinguishing primes are $113,181,197,281,293,307$, and $2642$ of the $17{,}981$ eligible primes up to $2\cdot10^5$
disagree:
\begin{center}
\artifact{grh_cubic_matching_and_precision_2026-11-27/results/cubic_checks.json}
\end{center}
\end{example}

\paragraph{Descent in the size and the rank.}
For nonzero balanced integral $c$ on $X=\F_\ell^r$ let $\nu(c)=\abs S/(\ell-1)$ be the number of projective directions on
which $H$ is nonzero (\cref{rec:lem:otrace}), and let $N_\ell(r,k)$ be the least $\nu(c)$ over such $c$ with $m\le k$, so that
$d_\ell(r,k)=(\ell-1)N_\ell(r,k)/\ell^r$; $N_\ell(r,\cdot)$ is nonincreasing, by padding.

\begin{lemma}[A balanced slice keeping a Fourier direction; \stU]\label{rec:lem:oslice}
Let $2\le s\le r$, and let $c\ne0$ satisfy $\nu(c)\le\ell^{s-1}$. There is an $s$-dimensional subspace $U\le X$ such that on every
coset $x+U$ the function $c_x(u)=c_{x+u}$ is nonzero, balanced, and has $\nu(c_x)\le\nu(c)$. Consequently some coset has
$1\le m(c_x)\le\floor{m/\ell^{r-s}}$.
\end{lemma}

\begin{proof}
Let $v_1,\dots,v_n$ represent the supported directions and draw $s-1$ vectors independently and uniformly from
$\ker v_1\subset X$. For $j>1$ the restriction of $v_j$ to $\ker v_1$ is a nonzero functional, and all draws lie in its
kernel with probability $\ell^{-(s-1)}$; the union of these events has probability at most $(n-1)/\ell^{s-1}<1$. Extend
the span of a successful tuple to an $(s-1)$-plane $W\subset\ker v_1$ and adjoin $y$ with $v_1(y)\ne0$ to obtain $U$. Every
supported direction restricts nontrivially to $U$, $v_1|_U$ vanishes on $W$ while no other supported direction
does, and $v_1(y)\ne0$; so $v_1$ is the only supported frequency restricting to $v_1|_U$. Fourier inversion on $X$ and
orthogonality on $U$ give, for $\lambda\in U^*$,
\[
 \sum_{u\in U}c_x(u)\zeta_\ell^{\lambda(u)}=\ell^{s-r}\sum_{v|_U=\lambda}H(v)\zeta_\ell^{-v(x)}.
\]
The coefficient at $\lambda=0$ vanishes, the one at $v_1|_U$ is a single nonzero term, for every coset, and all other
supported frequencies of $c_x$ are restrictions of supported frequencies of $c$. The positive masses of the
$\ell^{r-s}$ restrictions sum to $m$.
\end{proof}

\begin{theorem}[Leading-digit descent; \stU]\label{rec:thm:odescent}
For $a=\floor{\log_\ell k}$, $b=\floor{k/\ell^a}$ and $r\ge a+2$,
\[
 N_\ell(r,k)=N_\ell(r-a,b),\qquad d_\ell(r,k)=\ell^{-a}d_\ell(r-a,b).
\]
\end{theorem}

\begin{proof}
Put $s=r-a\ge2$. The function $\mathbf 1_{\set0}-\mathbf 1_{\set{e_1}}$ on $\F_\ell^s$ has exactly $\ell^{s-1}$ supported directions, those with
nonzero first coordinate. Pulling back under a quotient map $X\to\F_\ell^s$ multiplies the mass by $\ell^a$ and keeps the
number of directions, so $N_\ell(r,k)\le N_\ell(s,b)\le\ell^{s-1}$. Apply \cref{rec:lem:oslice} to a minimizer on $X$: a slice of
positive mass at most $\floor{k/\ell^a}=b$ has at least $N_\ell(s,b)$ and at most $N_\ell(r,k)$ directions.
\end{proof}

\begin{corollary}[Inadmissible digits; \stU]\label{rec:cor:ogap}
Let $a\le r-2$. Every value of $d_\ell(s,b)$, $s\ge2$, $b<\ell$, gives $d_\ell(r,k)$ for all sizes with leading digit $b$ and all
ranks $r\ge a+2$. If $b$ divides neither $\ell-1$ nor $\ell+1$ and $b_0<b$ is the largest integer dividing $\ell-1$ or $\ell+1$, then
\[
 \frac{(\ell-1)(\ell+2-b)}{\ell^{a+2}}\le d_\ell(r,k)\le\frac{(\ell-1)(\ell+1-b_0)}{\ell^{a+2}},
\]
with equality on the left when $b_0=b-1$.
\end{corollary}

\begin{proof}
Apply \cref{rec:thm:odescent} and \cref{rec:thm:osmall}(iv) in rank $r-a$ (an inadmissible digit is at least three); the
upper bound is the lift of the construction for $b_0$, padded.
\end{proof}

\paragraph{Small packets.}
For $1\le k<\ell$ let $B_\ell(k)$ be the largest integer at most $k$ that divides $\ell-1$ or $\ell+1$.

\begin{theorem}[The conic range; \stU]\label{rec:thm:oconic}
If $k<\ell$, $B=B_\ell(k)\ge3$ and $2k<3B$, then $d_\ell(r,k)=(\ell-1)(\ell+1-B)/\ell^2$ for every $r\ge2$.
\end{theorem}

\begin{proof}
Let $c\ne0$ be balanced with $m\le k$ and support $Z$; its weights are nonzero modulo $\ell$, as $\abs{c_x}\le k<\ell$. For $1\le j\le B$
put $P_j(v)=\sum_xc_xv(x)^j$, a form of degree $j$ over $\F_\ell$ on $V$. At an agreeing $v$ all projection fibres have signed
sum zero, so all $P_j$ vanish there. If some $P_j$ is nonzero, \cref{rec:lem:oproj} gives at least $(\ell+1-j)\ell^{r-2}\ge(\ell+1-B)\ell^{r-2}$
disagreeing directions, as required.

Suppose all $P_j$ vanish identically. Comparing coefficients (the multinomial coefficients in degrees below $\ell$
are units, and the constant term vanishes by balance) gives $\sum_xc_xf(x)=0$ for every polynomial $f$ of degree at
most $B$, also over $\overline{\F}_\ell$. So $Z$ has the Cayley--Bacharach property in degree $B$ (no polynomial of degree at most
$B$ separates a point of $Z$ from the others), and $\abs Z\le2k<3B$. By \cite[Theorem~3.8, Proposition~8.2]{Couvreur2012Dual}, a
dependent set of fewer than $3B$ points contains $B+2$ collinear points or $2B+2$ points on a plane conic. In the
conic case all of $Z$ lies on the conic: a point outside it is isolated by a form of degree at most two
vanishing on the conic and nonzero at the point, times one linear factor for each of the at most $B-3$ other
outside points (a point outside the plane is first separated from the plane by a linear form). In the line
case, a hyperplane over $\overline{\F}_\ell$ containing the line and missing the other points of $Z$ turns the relation into a
full-support relation of degree $B-1$ on the residual set $R$, with $\abs R\le2B-3$; if $R\ne\varnothing$,
\cite[Proposition~7.3]{Couvreur2012Dual} in degree $B-1$ gives $B+1$ collinear points of $R$, and the same separation
argument puts all of $R$ on that line. So $Z$ lies on a line, on a plane conic, or on two skew lines. Lines
spanned by two rational points are rational; in the conic case the plane is rational, a smooth conic is
rational because its Frobenius conjugate shares more than four rational points with it, and in the reducible
case both components are rational, each containing at least two rational points.

Affine changes preserve all matching proportions, and the restriction of a uniform functional to the span
of the translated support is uniform. If $Z$ is collinear only the zero functional on its direction can
agree, giving agreement $1/\ell$. For two skew lines $(t,0,0)$ and $(0,t,1)$ with signed arrays $f,g$, both are nonconstant
(a nonzero constant on all $\ell$ points uses mass $\ell$ of one sign); for a functional $(u,v,h)$ with $uv\ne0$ the
projections are a dilation of $f$ and an $h$-shift of a dilation of $g$, and at most one shift cancels them, since a
nonconstant array on $\Z/\ell\Z$ has no nonzero period; if exactly one of $u,v$ vanishes one array collapses to one bin
and the other does not; allowing all $\ell$ shifts when $u=v=0$, the agreement is at most $((\ell-1)^2+\ell)/\ell^3<1/\ell$.
Both are below $(1+(\ell-1)B)/\ell^2$.

It remains to show that a noncollinear signed array on a plane conic has at most $B$ matching directions.
A smooth affine conic with a rational point is, by an affine change, the parabola $(t,t^2)$, the split torus
$(t,t^{-1})$ or the norm-one torus $U=\set{t\in\F_{\ell^2}^\times:t^{\ell+1}=1}$. On the parabola a projection $ut+vt^2$ with $v\ne0$ has as
fibres the orbits of the reflection $t\mapsto-u/v-t$ ($v=0$ is injective); agreement means that the weight changes
sign under the reflection, and two such reflections would give a nonzero additive period, so there is at
most one matching direction. On the split torus a functional $ut+v/t$, $uv\ne0$, has as fibres the orbits of
$t\mapsto h/t$, $h=v/u$, and the axis functionals are injective; on the nonsplit torus every nonzero functional is
$t\mapsto\Tr(zt)$, with fibres the orbits of $t\mapsto z^{\ell-1}/t$, and $[z]\mapsto z^{\ell-1}$ is a bijection from directions onto $U$. In
either case the parameters $h$ for which the weight changes sign under $t\mapsto h/t$ form a coset of its
multiplicative stabilizer or are empty; the stabilizer acts freely and preserves the weight, so its order
divides $m$ and the group order $\ell\mp1$, and is at most $B$. For two intersecting lines, taken as axes, a matching
direction has both coefficients nonzero and forces the weight at the origin to vanish; agreement means that
one array on $\F_\ell^\times$ is minus a multiplicative translate of the other, and the working translates are empty or a
subgroup acting freely and preserving both weights, of order dividing $m$, hence at most $B$. For two parallel
lines $(t,0),(t,1)$, directions with nonzero first coordinate give translated arrays up to a common dilation, at
most one works, and with the remaining direction there are at most two. So in the plane at most $1+B(\ell-1)$
vectors agree, which is the bound.

For attainment, if $B\mid\ell-1$ take the two coordinate-axis copies of the subgroup of order $B$ of $\F_\ell^\times$ with
opposite signs, and if $B\mid\ell+1$ two distinct cosets of the subgroup of order $B$ of $U$ (\cref{rec:con:opack}); embed in
rank $r$ and pad to size $k$.
\end{proof}

For $\ell\ge7$ and $(\ell+1)/2\le k\le\floor{(3\ell+2)/4}$, $B_\ell(k)=(\ell+1)/2$ and the largest agreement is $(\ell^2+1)/(2\ell^2)$ in every rank
$r\ge2$; with \cref{rec:thm:odescent}, $d_{11}(r,k)=60/11^{a+2}$ for $8\cdot11^a\le k<9\cdot11^a$, $r\ge a+2$. At $a=0$ a witness lies on
$x^2+y^2=1$ modulo $11$: positive points $(\pm1,0),(\pm5,\pm3)$, negative points $(0,\pm1),(\pm3,\pm5)$ (independent signs), matching
functionals $(1,v)$ with $v\in\set{1,3,4,7,8,10}$, and two common entries on each side.

\begin{theorem}[Exponent eleven, size nine; \stU]\label{rec:thm:oeleven}
For every $r\ge2$, $d_{11}(r,9)=60/121$.
\end{theorem}

\begin{proof}
A padded six-point torus construction attains the value. With $B=6$, unless all moment forms of degrees one
to six vanish, \cref{rec:lem:oproj} gives at least $6\cdot11^{r-2}$ disagreeing directions. Otherwise $Z$ carries a full-support
sextic relation. If $\abs Z<18$ the proof of \cref{rec:thm:oconic} applies verbatim. Let $\abs Z=18$; then the reduced
packets are disjoint nine-sets and $c=\pm1$ on $Z$. The line and conic configurations still give at most six
planar matching directions, since no divisor of $10$ or $12$ lies between seven and nine. If eight points are
collinear, the at most ten residual points are Cayley--Bacharach in degree five, so by
\cite[Proposition~7.3]{Couvreur2012Dual} they contain seven collinear points if nonempty, and the linear-factor separator
puts them on that line. If fourteen points lie on a plane conic, a quadratic times at most three linear
factors separates any of the at most four outside points.

So assume no eight collinear and no fourteen on a plane conic. By \cite[Theorem~3.8(iii), Proposition~9.1]{Couvreur2012Dual}
(the set is linked and every proper subset is independent by \cite[Proposition~8.2]{Couvreur2012Dual}), $Z$ is the complete
intersection of a plane cubic $F$ and a sextic $G$ without common component; its plane is rational. Restricting
a cubic through $Z$ to a component $D$ of $F$, B\'ezout gives $6\deg D$ distinct points of $Z$ on $D$, none on two
components, while a cubic not containing $D$ meets it in at most $3\deg D$ points; so every cubic through $Z$
contains $F$, Frobenius fixes $F$ up to a scalar, and $F$ may be taken rational and reduced. A geometrically
irreducible singular cubic has a unique, hence rational, singular point, and each of the twelve rational lines
through it has at most one further rational point, too few for eighteen. If $F$ is smooth, Hasse's bound gives
$\#F(\F_{11})\le\floor{12+2\sqrt{11}}=18$, so the rational points are exactly $Z$. At a matching direction the occupied
parallel lines carry equal numbers of each sign and at most three points of $Z$, so there are nine occupied
lines with two points each; the third intersection with $F$ is rational, and cannot be a new point of $Z$, so
all nine lines are tangent to $F$. Their common point at infinity is off $F$, their nine contact points are
distinct and lie on the polar conic of that point, which is nonzero (in characteristic eleven a vanishing
directional derivative would make $F$ geometrically reducible); B\'ezout allows at most six contact points, so a
smooth $F$ gives no matching planar direction.

If $F$ is geometrically reducible without a rational line component, Frobenius cycles three lines and every
rational point lies on all three, which allows at most one point. So $F$ has a rational line component $L$
carrying exactly six points of $Z$, with the other twelve on the residual quadratic and off $L$. With an affine
parameter $t$ on $L$, the six parameters $S\subset\F_{11}$, $P_S(T)=\prod_{x\in S}(T-x)$, and $Q(t)$ the residual quadratic restricted
to $L$ (of degree at most two and nonzero on $S$), multiplication by the residual quadratic extends every
polynomial of degree at most four in $t$ to a sextic vanishing at the other twelve points, so
$\sum_{t\in S}c(t)Q(t)t^j=0$ for $0\le j\le4$. The kernel of this rank-five system is spanned by $(1/P_S'(t))_{t\in S}$, by Lagrange
interpolation, and as $c(t)^2=1$ the function $(Q(t)P_S'(t))^2$ is constant on $S$. Affine changes of $t$ and scalar
multiples of $Q$ preserve this condition, and completing the square (with $2$ a nonsquare modulo eleven) reduces
$Q$ to $1,t,t^2,t^2-1,t^2-2$. An exhaustive test of the $462,210,210,84,462$ six-subsets of the respective nonroots finds no
subset with constant $(Q(t)P_S'(t))^2$:
\begin{center}
\artifact{gram-euler-and-zero-filters/euler/check_euler.json}
\end{center}
These cases exhaust the cubic configurations.
\end{proof}

\begin{corollary}[Exponents seven and eleven; \stU]\label{rec:cor:oseveneleven}
For $\ell\in\set{7,11}$ and $B_\ell(b)$ the largest divisor of $\ell-1$ or $\ell+1$ at most $b$,
\[
 d_\ell(r,k)=\frac{(\ell-1)(\ell+1-B_\ell(b))}{\ell^{a+2}}\quad(r\ge a+2),\qquad d_\ell(r,k)=\frac{\ell-1}{\ell^r}\quad(r\le a+1).
\]
\end{corollary}

\begin{proof}
For $\ell=7$ the only inadmissible digit is $5$ (\cref{rec:cor:ogap}). For $\ell=11$ the inadmissible digits are $7,8,9$ with
$B_{11}=6$: \cref{rec:thm:oconic} covers $7$ and $8$, and \cref{rec:thm:oeleven} covers $9$, each lifted by \cref{rec:thm:odescent}.
Admissible digits are \cref{rec:thm:odigit}(ii), and saturated ranks \cref{rec:thm:psat}.
\end{proof}

In unsaturated rank the minimum is $\ell^{-a-2}$ times the following entry, and its complement is the largest
agreement:
\[
\begin{array}{c|rrrrrrrrrr}
 & \multicolumn{10}{c}{\text{leading digit }b}\\
 \ell &1&2&3&4&5&6&7&8&9&10\\\hline
 7&42&36&30&24&24&12&-&-&-&-\\
 11&110&100&90&80&70&60&60&60&60&20
\end{array}
\]
For instance, at exponent eleven two products of nine $L$-functions agree at most at $61/121$ of the primes in
every rank at least two, and for sizes $88$ to $98$ at most at $1-60/1331$ in every rank at least three. An
exhaustive enumeration of the $6{,}674{,}031$ pairs of size three at $(\ell,r)=(3,3)$ and explicit conic witnesses at
eight primes are in \artifact{gram-euler-and-zero-filters/euler/check_euler.json}.

\paragraph{The second level.}
Normalize by the reduced mass. By \cref{rec:thm:pmass} with $X[\ell]=X$ (\cref{rec:lem:oradon}), $\delta/\beta_m\ge\norm c_\infty$ for every
distinct pair, with $\delta=\beta_m$ exactly for the \emph{coset pairs} $c=\mathbf 1_{x+P}-\mathbf 1_{y+P}$, $P$ a subspace and $y-x\notin P$ (so
$m=\abs P$); a repeated coefficient already gives ratio at least $2$. For $c=\mathbf 1_A-\mathbf 1_B$ with disjoint $m$-sets let $P=\set{t:c(x-t)=c(x)\ \forall x}$ be its period subspace and
$h=m/\abs P$ the number of $P$-cosets on each side. Passing to $X/P$ divides $m$ by $\abs P$, multiplies $\delta$ by $\abs P$ and
keeps the ratio. Put $T_\ell=2(1-1/\ell)$.

\begin{theorem}[Isolation of the two lowest levels; \stU]\label{rec:thm:olevel}
A simple pair with $h\ge2$ has $\delta/\beta_m\ge(3h-2)/(2h-1)$. For $r\ge2$ the least value of $\delta/\beta_m$ over
all pairs other than the coset pairs and all sizes is $T_\ell$, and it is attained exactly, after passing to the quotient by the period
subspace and up to translation and exchange of sides, by the opposite vertex pairs $\set{0,u+w}$, $\set{u,w}$ of a
parallelogram ($u,w$ independent), and by two intersecting affine lines each with the common point deleted.
At $\ell=3$ the two descriptions coincide and $T_3=4/3$. In rank one $\delta/\beta_m=m$.
\end{theorem}

\begin{proof}
Fix $a\in A$ and $g_a=J_a/m$. On $S$, dropping the contribution of $B$,
\begin{equation}\label{rec:eq:oanchor}
 \mathbf 1_S-g_a\ge\frac1m\sum_{y\in A}\mathbf 1_{S\cap\set{v(y-a)\ne0}},
\end{equation}
and off $S$ both sides vanish. For $y\notin a+P$ and $t=y-a$, $c_t(x)=c(x)-c(x-t)$ is nonzero and balanced with
transform $H(v)(1-\zeta_\ell^{v(t)})$, so its disagreement set is the summand; its positive part comes from $A\cup(B+t)$
and the $\abs P$ points of $y+P$ cancel against $A+t$, so its mass is $m_t\le2m-\abs P$. By \cref{rec:lem:oradon} each
summand has mean at least $(1-1/\ell)/(2m-\abs P)$, and the $m-\abs P$ such $y$ give
$\delta/\beta_m\ge1+(m-\abs P)/(2m-\abs P)=(3h-2)/(2h-1)$; the terms with $y\in a+P$ vanish since $S\subset P^\perp$.

Assume $P=0$ after the quotient. If $c_t=\mathbf 1_{u'+D}-\mathbf 1_{v'+D}$ with $\abs D=N$, then
\begin{equation}\label{rec:eq:oflat}
 N\le h/(1-1/\ell).
\end{equation}
Indeed $c_t$ sums to zero along every line parallel to $t$, so $t\notin D$, both cosets project onto the same $N$ points
modulo $\langle t\rangle$, and $v'+D=u'+D+st$ for one $1\le s\le\ell-1$; on each of these $N$ lines $c$ takes values $0,1$ on
$s$ points or $-1,0$ on $\ell-s$ points, so with $j$ lines of the first type
$h\ge\max\set{js,(N-j)(\ell-s)}\ge Ns(\ell-s)/\ell\ge N(1-1/\ell)$.

Let $\alpha$ be the infimum of $\delta/\beta_m$ over pairs not of coset form, at fixed $\ell$ and all ranks. The parallelogram has
$H=(1-\zeta_\ell^{v(u)})(1-\zeta_\ell^{v(w)})$ up to a unit, $\delta=(1-1/\ell)^2$ and $m=2$, so $\alpha\le T_\ell$. Repeated coefficients give
ratio at least $2$, a quotient of rank one gives ratio $h\ge2$, and for quotient rank at least two and
$2\le h\le\ell-1$, \cref{rec:thm:osmall} gives $\delta/\beta_h\ge h(\ell+1-h)/\ell\ge T_\ell$, with equality only for $h\in\set{2,\ell-1}$. For
$h\ge\ell$ average \eqref{rec:eq:oanchor}: a summand with $c_t$ not of coset form has mean at least $\alpha\beta_N$ with $N\le2h-1$,
and one of coset form has mean $\beta_N$ with \eqref{rec:eq:oflat}, so
\[
 \delta/\beta_h\ge\min\set{1+(1-1/h)(1-1/\ell),\ 1+\alpha(h-1)/(2h-1)}.
\]
The first entry exceeds $T_\ell$ by at least $1/\ell^2$; if $\alpha\le T_\ell$ the second exceeds $\alpha$ by at least
$1-T_\ell\ell/(2\ell-1)=1/(2\ell-1)$. Hence every ratio is at least $\min\set{T_\ell,\alpha+1/(2\ell-1)}$, so $\alpha=T_\ell$, and $h\ge\ell$ never
attains it. At $h=2$ the vector sums must agree, since otherwise $\delta\ge1-1/\ell$; after translation the pairs are
$\set{0,u+w}$, $\set{u,w}$, with $u,w$ independent as proportional ones give $1-1/\ell$. At $h=\ell-1\ge4$, $\delta=D_1(\ell-1)$, so
\cref{rec:thm:osmall}(ii) puts the $2\ell-2>\ell+1$ points in a plane on a reducible conic with monochromatic lines;
parallel lines would need an additive subgroup of order $\ell-1$, so the lines intersect, the common point lies in
neither packet, and each packet is its line minus that point. Conversely $c$ is then the difference of the two
full line indicators, with Fourier support the two punctured annihilator lines, $\delta=2(\ell-1)/\ell^2$ and ratio $T_\ell$.
\end{proof}

\paragraph{Factors of one primitive conductor.}

\begin{theorem}[Primitive common-conductor floor; \stU]\label{rec:thm:oprim}
If every character of both packets is primitive modulo the same $Q$, then $\delta(A,C)\ge(\ell-1)2^{r-1}/\ell^r$. For every
$Q$ supporting a primitive character of order $\ell$ this is the exact minimum in the full exponent-$\ell$ group
modulo $Q$, of rank $r=\omega(Q)$, attained for every $k\ge2^{r-1}\ell^{r-1}$. In local coordinates, an extremal
disagreement set is, in rank two, the union of two punctured lines of nonzero finite slopes, and in rank $r\ge3$
the set $\set{(t,\epsilon_2d_2t,\dots,\epsilon_rd_rt):t\in\F_\ell^\times,\ \epsilon_i=\pm1}$ for some $d_2,\dots,d_r\in\F_\ell^\times$.
\end{theorem}

\begin{proof}
For odd $\ell$ the order-$\ell$ characters of each prime-power factor of $Q$ form a group of dimension at most one,
and the restrictions separate $X$; take $r$ independent ones as coordinates. A primitive character is
nontrivial at every factor, so $c$ vanishes when a coordinate is zero, and orthogonality gives
$\sum_{t\in\F_\ell}H(v+te_i)=0$ for all $i$ and $v$. A nonzero array on $\F_\ell^n$ with zero sums along all coordinate lines has at
least $2^n$ nonzero entries (two nonzero slices in any coordinate, and induction); a nonzero slice at $v_1\ne0$
exists, Galois conjugation gives all $\ell-1$ such slices, and each has at least $2^{r-1}$ nonzero entries.
For sharpness, $Q$ is a product of distinct primes $q\equiv1\pmod\ell$, possibly times $\ell^2$, and at each factor the
nontrivial local order-$\ell$ characters are exactly the primitive ones. Put $\lambda(t)=1$ for even $r$,
$\lambda(t)=\zeta_\ell^t-\zeta_\ell^{-t}$ for odd $r$, $J(0,v')=0$ and $J(t,v')=\lambda(t)\prod_{i=2}^r(\mathbf 1_{v_i=t}-\mathbf 1_{v_i=-t})$ for $t\ne0$; the line
sums vanish, using $\lambda(-t)=(-1)^r\lambda(t)$. With $u_\epsilon=x_1+\sum_{i\ge2}\epsilon_ix_i$ and $e_\epsilon=\prod_{i\ge2}\epsilon_i$, the inverse transform
is $c_x=\sum_\epsilon e_\epsilon\mathbf 1_{u_\epsilon=0}$ for even $r$ and $\sum_\epsilon e_\epsilon(\mathbf 1_{u_\epsilon=1}-\mathbf 1_{u_\epsilon=-1})$ for odd $r$: integers of sum zero,
masses at most $2^{r-1}\ell^{r-1}$, vanishing at every character with a zero coordinate, with exactly $(\ell-1)2^{r-1}$
disagreement points; pad with a primitive character. For the shape, equality in the array bound forces
exactly two nonzero, pointwise opposite slices in each coordinate, so the support of the slice $v_1=1$ is a
product $D_2\times\dots\times D_r$ of two-element subsets of $\F_\ell^\times$, and the slice $v_1=t$ has support $tD_2\times\dots\times tD_r$. For
$r\ge3$ the line through a chosen point of the slice $v_1=1$ in the first direction has a second support point
$t\ne1$, so each chosen element divided by $t$ is its partner; the common ratio squares to one and is not one,
so $D_i=\set{\pm d_i}$.
\end{proof}

\begin{theorem}[Primitive cubic packets; \stU]\label{rec:thm:ocubic}
Let $\ell=3$ and $k\ge1$, and put $d=\min\set{r,\floor{\log_2(2k)}}$. Two distinct size-$k$ packets of primitive
characters of one conductor, in a rank-$r$ space, satisfy $\delta\ge(2/3)^d$, with equality achieved for every $k$ in the full cubic
group of every conductor supporting a primitive cubic character. The common-conductor minimum
$(2/3)^r$ is first reached at $k=2^{r-1}$.
\end{theorem}

\begin{proof}
In local coordinates every $x$ with $c_x\ne0$ has $x_i\in\set{1,2}$; with $z_i=\zeta_3^{v_i}$, $H(v)/\prod_iz_i$ is a multilinear
polynomial with $n\le2k$ monomials on the cube of cube roots of unity. Such a nonzero polynomial is nonzero
on a fraction at least $(2/3)^{\min(r,\floor{\log_2n})}$: split $P=P_0+z_iP_1$ in a coordinate where both parts are nonzero, apply
induction to the part with at most $\floor{n/2}$ monomials, and note that where it is nonzero at most one value of
$z_i$ kills $P$. For equality take the coefficients of $\prod_iz_i\prod_{i\le d}(z_i-1)$: disjoint packets of size $2^{d-1}$ of
primitive characters, disagreeing exactly where $z_i\ne1$ for $i\le d$; pad.
\end{proof}

In rank three the primitive cubic minima for $k=1,2,3,4$ are $2/3,4/9,4/9,8/27$, against $2/3,4/9,2/9,2/9$
without primitivity.

\begin{proposition}[Two factors; \stU]\label{rec:prop:otwo}
The agreement densities of distinct pairs of size two are $1/\ell$ in rank one; $1/\ell$, $(2\ell-1)/\ell^2$, $1/\ell^2$ in rank
two; and $1/\ell$, $(2\ell-1)/\ell^2$, $1/\ell^2$, $(2\ell-1)/\ell^3$ in rank at least three.
\end{proposition}

\begin{proof}
Write the packets as $\set{s\pm u}$ and $\set{t\pm w}$; as $2$ is invertible they agree exactly on
$(s-t)^\perp\cap\bigl((u-w)^\perp\cup(u+w)^\perp\bigr)$. Inside the first hyperplane (or the whole space if $s=t$) the conditional
fraction is $1$, $1/\ell$ or $(2\ell-1)/\ell^2$ according as one of the other forms vanishes there, their restrictions are
proportional, or independent; the three vectors can be prescribed freely.
\end{proof}

\begin{example}[Fifteen quintic factors]\label{rec:ex:oquintic}
Let $Q=11\cdot31\cdot41=13981$, let $\chi_1,\chi_2,\chi_3$ send the least primitive roots $2,3,6$ of the three primes to
$e^{2\pi i/5}$, and write $(i,j,t)$ for $\chi_1^i\chi_2^j\chi_3^t$. The packets $A=\set{(1,0),(2,2),(2,3)}\times\F_5$ and $C=-A$ consist of
characters of exact order five (the first two coordinates are the order-three subgroup of the norm-one
conic in $\F_5[\theta]/(\theta^2-2)$ and its negative). Exactly $12$ of the $125$ evaluation vectors distinguish them, so the two
products of fifteen $L$-functions have the same Euler factor at $113/125$ of the primes, the maximum in rank
three by \cref{rec:thm:odigit} ($a=1$, $b=3\mid6$). The smallest distinguishing primes are $3,79,163,173,191$, and $906$ of the
$9589$ unramified primes up to $10^5$ distinguish them.
\end{example}

\begin{example}[Five primitive septic factors]\label{rec:ex:oseptic}
For $Q=29\cdot43=1247$ and coordinates given by order-seven characters at the two primes, the packets
$\set{(4,3),(2,3),(3,4),(3,2),(3,3)}$ and $\set{(5,5),(1,5),(5,1),(1,1),(3,3)}$ consist of primitive characters of
conductor $1247$. After translating by $-(3,3)$ and cancelling the common entry, they are $\set{\pm1,\pm i}$ and
$(2+2i)\set{\pm1,\pm i}$ in $\F_7[i]$, where $2+2i$ has norm one; four of the eight directions agree, and the products
agree at $25/49$ of the primes, the maximum by \cref{rec:thm:osmall}(iv) (\cref{rec:cor:oseven}). For the translated four-element packets the
primes up to $10^6$, apart from $29,43$, give $38{,}440$ disagreements and $40{,}056$ agreements.
\end{example}

Exhaustive enumerations confirm the minima on all $220{,}813{,}980$ pairs in ten small parameter sets and
$7{,}913{,}653$ primitive cubic pairs, all $189{,}471$ pairs in further small planes with $103$ constructions, and all
$4{,}276{,}350$ pairs at $(\ell,r,k)=(5,2,3)$:
\begin{center}
\artifact{odd-order-euler-factor-matching/checks.json}\\
\artifact{euler-packet-rigidity/checks/output.txt}\\
\artifact{odd_prime_order_euler_packets/verification.txt}
\end{center}
Independent exact enumerations, including $17{,}735{,}976$ cubic pairs through rank four with $5285$ extremal
structures, the absence of a planar five-point configuration at $\ell=7$ ($5280$ permutation pairs), and the
split, nonsplit and line constructions for primes through $31$, are in
\begin{center}
\artifact{grh_cubic_matching_and_precision_2026-11-27/results/cubic_checks.json}\\
\artifact{prime_exponent_euler_agreement/verification.txt}\\
\artifact{odd_prime_euler_agreement/verification.json}\\
\artifact{cubic_quintic_euler_agreement/verification.txt}
\end{center}
Related matching problems for Galois representations are studied in \cite{Walji2016Matching}; the planar bound
for $k<\ell$ also follows from Lev's bound on perfect directions \cite{Lev2020Point}, and the projective
counts are the minimum distance of projective Reed--Muller codes \cite{GhorpadeLudhani2024Minimum}.

Let $w$ and $M_Y$ be as in \eqref{rec:eq:cgram}. At exponent $\ell$, \cref{rec:prop:pwindow} gives
$M_Y(\mathbf 1_S)\ge\norm c_\infty\beta_mY-2C_w\sqrt Y\log(2Q)$ for every distinct pair, and $M_Y(\mathbf 1_S)=\beta_mY+O(C_w\sqrt Y\log(2Q))$ for the coset
pairs. The following theorem adds the gains of the higher levels.

\begin{theorem}[Distinguishing primes in one window; \stG]\label{rec:thm:ofinite}
Let $w$ and $M_Y$ be as in \eqref{rec:eq:cgram}, and let all implied constants depend only on $w$ (in particular not
on $\ell,r,k,m,Q$). For distinct packets:
\begin{enumerate}
\item $M_Y(\mathbf 1_S)\ge\Lambda\beta_mY-O(\sqrt Y\log(2Q))$ with $\Lambda=(3h-2)/(2h-1)\ge4/3$ for simple pairs not of coset form, $h$ as in
\cref{rec:thm:olevel};
\item $M_Y(\mathbf 1_S)\ge(\ell-1)\ell^{-r}Y-O(\sqrt Y\log(2Q))$, and for $r\ge2$, $k\le\ell$, $M_Y(\mathbf 1_S)\ge D_1(k)Y-O(k\sqrt Y\log(2Q))$;
\item for pairs not of coset form, $M_Y(\mathbf 1_S)\ge T_\ell\beta_mY-O(\log(2\ell)\sqrt Y\log(2Q))$, and two-sided
$M_Y(\mathbf 1_S)=T_\ell\beta_mY+O(\sqrt Y\log(2Q))$ for the second-level pairs of \cref{rec:thm:olevel};
\item each interval $(Y/2,Y)$ has a distinguishing prime once $Y\ge C\min(m,\ell^{r-1})^2\log^2(2Q)$;
\item a pair attaining $D_\ell(r,k)$ with $a\le r-2$ has $\norm{\mathbf 1_S}_A=2b(\ell+1-b)(\ell-1)/\ell^2$ exactly, so
$M_Y(\mathbf 1_S)=Y\delta+O(2b\ell^a\delta\sqrt Y\log(2Q))$; at saturation the norm is $2(1-1/\ell)$. Given only an upper bound
$K_0$ on the reduced size and the promise of an unsaturated equality case, rounding $M_Y(\mathbf 1_S)/Y$ to the nearest
admissible value of $D_\ell$ recovers $\delta$ and the reduced size $b\ell^a$ once $Y\ge C_\ell K_0^2\log^2(2Q)$, uniformly in the rank;
\item with prime weights $(\log p)w(p/Y)/M_Y(1)$ and weighted Hamming distance, every word of local factors
within weighted distance below $(1-\eta)D_0/2$ of a packet, $D_0=(\ell-1)/(\ell\min(k,\ell^{r-1}))$, determines it once
$Y\ge C\eta^{-2}\min(k,\ell^{r-1})^2\log^2(2Q)$; more generally erasures of weight $u$ and errors of weight $e$ are tolerated when
$u+2e<(1-\eta)D_0$. Where $d_\ell(r,k)=D$ is attained by a coset or saturated pair or a small equality
case, no uniform correction radius above $D/2$ survives as $Y\to\infty$.
\end{enumerate}
\end{theorem}

\begin{proof}
By \eqref{rec:eq:cgram} and $w\ge0$, a pointwise minorant $f\le\mathbf 1_S$ gives $M_Y(\mathbf 1_S)\ge Y\mathbb Ef-C_w\norm f_A\sqrt Y\log(2Q)$.
(i) Fix $a\in A$; $g_a=J_a/m$ is a minorant of mean $\beta_m$ and norm at most $2$ (\cref{rec:lem:oradon}). Add to $g_a$ the sum over $y\notin a+P$ of $1/m$ times the
minorant $g$ of the summand of \eqref{rec:eq:oanchor}, anchored and signed so that its largest coefficient is positive;
the result is below $\mathbf 1_S$, of norm at most $4$, with the mean of the proof of \cref{rec:thm:olevel}.
(ii) Use $\mathbf 1_\Lambda-\mathbf 1_{\set0}$ for a line $\Lambda$ with $\Lambda\setminus\set0\subset S$ (norm at most $2$), respectively
$(1-\mathbf 1_{H_1})(1-\sum_{j=2}^d\mathbf 1_{H_j})$ for the anchored hyperplanes of \cref{rec:thm:osmall}, which is $1$ off their union and
nonpositive on it, has mean $(\ell-1)(\ell+1-d)/\ell^2$ and norm at most $2d$.
(iii) Let $s=1-1/\ell$. With a repeated coefficient use \cref{rec:prop:pwindow}, as $\norm c_\infty\ge2>T_\ell$. Otherwise pass to the period quotient, with $n$ points of
each sign; extension by zero from $P^\perp$ divides the mean by $\abs P$ and keeps the norm. For $2\le n\le(\ell-1)/2$ the
minorant of (ii), scaled by $T_\ell\beta_n/D_1(n)\le1$, has norm at most $2T_\ell\ell/(\ell+1-n)\le8$. For larger $n$ iterate
\eqref{rec:eq:oanchor} starting from $L_0=1$: a summand of coset form has $m_t\le n/s$ by \eqref{rec:eq:oflat}, so $1+(n-1)/m_t\ge T_\ell$
(for $n\le\ell-2$ use $m_t\le n$), and its indicator has norm at most $2$; a summand not of coset form has, by the
previous stage, a minorant of mean at least $L_is/m_t$ with $m_t\le2n-1$. This gives minorants of mean
$L_{i+1}\beta_n$, $L_{i+1}=\min\set{T_\ell,1+\alpha L_i}$ with $\alpha=(\ell-1)/(2\ell)\le(n-1)/(2n-1)$, the norm growing by at most $2$ per step.
The fixed point $2\ell/(\ell+1)=T_\ell+2/(\ell(\ell+1))$ is approached with ratio below $\tfrac12$, so $\ceil{\log_2(\ell(\ell+1))}$ steps reach $T_\ell$.
For the two-sided claim, the parallelogram indicator is a product of two indicators $1-\mathbf 1_{\text{hyperplane}}$, and that
of two punctured lines is $\mathbf 1_\Lambda+\mathbf 1_{\Lambda'}-2\mathbf 1_{\set0}$; both have norm at most $4$, also after pullback and extension.
(iv) Combine \cref{rec:prop:pwindow} and (ii) with $\beta_m\ge2/(3m)$ and $(\ell-1)\ell^{-r}\ge\tfrac23\ell^{1-r}$.
(v) By \cref{rec:thm:odigit}(ii), $S$ is the inverse image of a union $\Pi$ of $d'=\ell+1-b$ punctured lines in a plane. The
coefficient of $\mathbf 1_\Pi$ is $d'(\ell-1)/\ell^2$ at $0$, $(\ell-d')/\ell^2$ at the $d'(\ell-1)$ nonzero characters trivial on a selected line and
$-d'/\ell^2$ at the other $b(\ell-1)$; the sum of absolute values is $2bd'(\ell-1)/\ell^2$, and pullback and extension by zero
preserve the norm, while $\delta=d'(\ell-1)/\ell^{a+2}$. By (ii) of \cref{rec:thm:odigit} the reduced size is $b\ell^a\le K_0$, so $a\le A_0=\floor{\log_\ell K_0}$.
The values $(\ell-1)(\ell+1-b)/\ell^{a+2}$ decrease strictly in the lexicographic order of $(a,b)$, consecutive ones differing by
$(\ell-1)/\ell^{a+2}$ also across $(a,\ell-1)\to(a+1,1)$, so the candidates with $a\le A_0$ are $(\ell-1)/\ell^{A_0+2}$-separated, while the norm is
at most $\ell$ and the error is $O(\ell\log(2Q)/\sqrt Y)$, below half the separation for $Y\ge C_\ell\ell^{2A_0}\log^2(2Q)$.
(vi) Two packets within weighted distance $e$ of one word are within $2e$ of each other, and two packets
agreeing with one word off erasures of weight $u$ and up to errors of weight $e$ are within $u+2e$; by (ii) and
\cref{rec:prop:pwindow} any two distinct packets are at relative distance at least $(D_0-2\eps)/(1+\eps)$, $\eps=C\log(2Q)/\sqrt Y$. For sharpness, on the
agreeing primes of an extremal pair use the common factor and split the differing primes into two parts of
nearly equal weight, assigning one packet to each; the largest single weight is $O(\log Y/Y)$, and the
distances of extremal pairs converge to $D$ by the two-sided bounds.
\end{proof}

In rank two a disagreement set is a union of at most $\ell+1$ punctured lines, so its indicator has norm at most
$2(\ell+1)$; all distances of a fixed rank-two size converge with error $O(\ell\log(2Q)/\sqrt Y)$, and $24/49$ is also the limiting
minimum of the windowed problem at $(\ell,r,k)=(7,2,5)$.

\subsection{Open questions}\label{eul:sec:open}
{\raggedright\sources{dirichlet_packet_reconstruction; euler-packet-rigidity; odd-order-euler-factor-matching; odd_prime_order_euler_packets; cubic_quintic_euler_agreement; odd_euler_agreement; odd_prime_euler_agreement; odd_prime_euler_agreement_2; odd_prime_euler_packets; prime_exponent_euler_agreement; prime_order_euler_agreement; grh_cubic_matching_and_precision_2026-11-27; euler_agreement_2026-11-27; gram-euler-and-zero-filters}}

\begin{enumerate}
\item \emph{Euler packets of prime order.} For Euler packets of odd prime order the least disagreement density and all extremal pairs are known
whenever the leading base-$\ell$ digit $b$ of the size divides $\ell\pm1$ or the size is saturated
(\cref{rec:thm:odigit,rec:thm:oplanar,rec:prop:psateq}), and the two lowest normalized levels by \cref{rec:thm:olevel}; for the other
digits \cref{rec:thm:odescent} reduces the minimum to sizes $k<\ell$ in rank at least two, \cref{rec:thm:oconic} settles
them when $2k<3B_\ell(k)$, and every rank and size is settled for $\ell\le11$ (\cref{rec:cor:oseveneleven}). What is the
minimum for $k<\ell$ when $2k\ge3B_\ell(k)$, where the support can reach the third configuration of Couvreur's
hierarchy, and which complete intersections carry a balanced integral relation with unit coefficients and
many matching directions? At an inadmissible digit, what are all minimizers, and when must a minimizer be a
pullback? Which minimum-weight words of the projective Reed--Muller code, or which second-weight words, are
reductions of integral spectra of a prescribed mass, and in rank two for which $(\ell,k)$ with $2k\ge3B_\ell(k)$ does the
minimum exceed the gap value of \cref{rec:thm:osmall}(iv)? What is the next normalized level above $2(1-1/\ell)$, the next
density at a fixed size $k=\ell^s$ (at least $T_\ell\beta_k$), and the whole spectrum, already the second-smallest density
and the range below twice the minimum for $\ell=3,5$? Is there a sharp primitive bound at each packet size for
$\ell\ge5$ refining \cref{rec:thm:oprim}? How large can the least distinguishing prime of an extremal pair be, can
the extremality promise in \cref{rec:thm:ofinite}(v) be removed while keeping a rank-free scale polynomial in
$k$, and is there an efficient decoder at the correction radius of \cref{rec:thm:ofinite} or
\cref{rec:cor:ccorrect}? At composite exponent the divisor-indexed traces of \cref{rec:prop:cdivisor} decide the
local polynomial, and the minimum is known for cyclic groups at every size (\cref{rec:thm:cyclic}) and for
abelian $\ell$-groups at the sizes $\ell^u$ (\cref{rec:thm:psat}). What is $d_X(k)$ for an arbitrary finite abelian $X$ and every $k$,
in particular when several primary components have rank above one? What is the spectrum of an abelian
$\ell$-group at sizes that are not powers of $\ell$, where the mass and saturation bounds of
\cref{rec:thm:pmass,rec:thm:psat} still hold? Is every cyclic extremizer a parity box?
\item \emph{Quadratic packets.} Can one density measurement in \cref{rec:thm:qmatchfinite} be made at a smaller scale?
\end{enumerate}
